\exercisesection{\S~5}

\begin{enumerate}
\def\labelenumi{\arabic{enumi})}
\item
  假定 G 是单连通的；G 的单房是 G 中形如 \(\exp(A)\) 的子集，其中 A 是 g 的某个嘉当子代数的一个单房。

\begin{enumerate}
\def\labelenumi{\alph{enumi})}
\item
  证明 \(\mathbf{G}\) 的诸单房构成 \(\mathrm{G}_r\) 的一个分划。
\item
  G 的每个单房都包含于唯一的极大环面。
\item
  证明 G 的包含于 T 的诸单房构成 \(T_{r}\) 的一个分划，且这些单房组成的集合是 W 下的一个主齐性空间。
\item
  G 的正则元素的每个共轭类在每个单房中都有唯一的点。
\item
  设 E 是 G 的一个单房。证明 E 是一个可缩空间；若 g 是单的，则 E 同胚于一个欧几里得空间中的开单纯形。
\end{enumerate}

\begin{enumerate}
\def\labelenumi{\arabic{enumi})}
\setcounter{enumi}{1}
\tightlist
\item
  \begin{enumerate}
  \def\labelenumii{\alph{enumii})}
  \tightlist
  \item
    设 E 是一个单连通拓扑空间，\(u : E \to G\) 是满足 \(u(E) \subset G_r\) 的连续映射。证明 u（《一般拓扑学》第 XI 章）同伦于以 e 为值的常值映射（考察覆盖 \(\varphi_r : (G/T) \times t_r \to G_r\)；把 u 提升为 \(\tilde{u} : E \to (G/T) \times t\)，再利用 t 可缩这一事实）。
  \end{enumerate}
\end{enumerate}

\(^{*}b)\) 证明群 \(\pi_{2}(G)\) 为零。

（设 \(u: \mathbf{S}_2 \to \mathbf{G}\) 是一个 \(\mathrm{C}^\infty\) 类映射；利用命题 1 与横截性定理，证明 \(u\) 同伦于一个像包含于 \(\mathrm{G}_r\) 中的映射，然后应用 \(a).)_*\)

\begin{enumerate}
\def\labelenumi{\arabic{enumi})}
\setcounter{enumi}{2}
\tightlist
\item
  设 \(f: \tilde{\mathrm{G}} \to \mathrm{G}\) 是 \(\mathrm{G}\) 的一个泛覆盖，\(\pi\) 是 \(f\) 的核（同构于 \(\pi_1(\mathrm{G})\)）；置 \(\mathrm{C} = \mathrm{C}(\mathrm{G})\) 与 \(\tilde{\mathrm{C}} = \mathrm{C}(\tilde{\mathrm{G}})\)。设 \(\sigma\) 是
\end{enumerate}

G 的一个自同构，\(\tilde{\sigma}\) 是 \(\tilde{G}\) 的一个自同构（由 \(\sigma\) 诱导）。用 \(G_{\sigma}, C_{\sigma}, \tilde{G}_{\sigma}, \tilde{C}_{\sigma}\) 分别表示 G, C, \(\tilde{G}\)，\(\tilde{C}\) 中被 \(\sigma, \sigma, \tilde{\sigma}, \tilde{\sigma}\) 固定的点组成的集合。

\begin{enumerate}
\def\labelenumi{\alph{enumi})}
\item
  证明 \((\mathrm{G}_{\sigma})_0\)（\(\mathrm{G}_{\sigma}\) 的单位元连通分支）是 \(f(\tilde{\mathrm{G}}_{\sigma})\)。
\item
  用 s 表示 Z-模 \(\pi\) 的由 \(\tilde{\sigma}\) 诱导的自同态。证明商群 \(\mathrm{G}_{\sigma}/(\mathrm{G}_{\sigma})_{0}\) 同构于 \(\operatorname{Coker}(1-s)\) 的一个子群，特别地它是交换的。若 1-s 是满射，则 \(G_{\sigma}\) 是连通的。
\item
  设 \(n \in N\) 满足 \(s^{n} = \operatorname{Id}_{\pi}\)。证明群 \(\mathrm{G}_{\sigma}/(\mathrm{G}_{\sigma})_{0}\) 可与商 \(\operatorname{Ker}(1 + s + \cdots + s^{n-1})/\operatorname{Im}(1 - s)\) 的一个子群等同；由此推出它被 n 零化。~若 n 与 \(\pi\) 的挠子群的阶互素，则 \(G_{\sigma}\) 是连通的。
\end{enumerate}

\(d\) ) 利用《代数学》（Algèbre）第 X 章第~194 页的习题 23，重新得出 \(b\)) 与 \(c\)) 的结果。

\begin{enumerate}
\def\labelenumi{\alph{enumi})}
\setcounter{enumi}{4}
\tightlist
\item
  证明在下列每种情形中 \(\mathrm{G}_{\sigma}\) 都是连通的：
\end{enumerate}

\begin{enumerate}
\def\labelenumi{(\roman{enumi})}
\item
  G 是 \(\mathrm{A}_{2n}\) 型（\(n \geq 1\)）的半单群且 \(\sigma\) 不是内的；
\item
  G 是 \(\mathrm{E}_6\) 型的半单群且 \(\sigma\) 不是内的；
\item
  G 是 \(\mathrm{D}_4\) 型的半单群且 \(\sigma\) 是一个三重性（即模 \(\operatorname{Int}(\mathbf{G})\) 后的阶为 3）。
\end{enumerate}

\(f\) ) 定义一个从 \(\mathrm{C}_{\sigma} / (\mathrm{C}_{\sigma}\cap (\mathrm{G}_{\sigma})_{0})\) 到 \(((1 - \tilde{\sigma})\tilde{\mathrm{C}}\cap \pi) / (1 - \tilde{\sigma})\pi .\) 的同构。由此推出，若 \(1 - s\) 不是满射且 \(\pi \subset (1 - \tilde{\sigma})\mathrm{C}\)，则 \(\mathrm{C}_{\sigma}\) 不包含于 \((\mathrm{G}_{\sigma})_{0}\)。对 \(\mathrm{G} = \mathbf{SO}(2n,\mathbf{R})\) 且 \(\sigma\) 非内时，有 \(-I_{2n}\notin (\mathrm{G}_{\sigma})_{0}\)，且 \(\mathrm{G}_{\sigma}\) 不是连通的。

\begin{enumerate}
\def\labelenumi{\arabic{enumi})}
\setcounter{enumi}{3}
\tightlist
\item
  假定 G 是半单的。设 e 是 G 的一个标架，\(\Phi\) 是 G 的一个尊重该标架的自同构群；用 H 表示 G 中由被 \(\Phi\) 固定的元素组成的子群。
\end{enumerate}

\begin{enumerate}
\def\labelenumi{\alph{enumi})}
\item
  证明 \(H_{0}\) 是半单的；若 G 是单连通的，则 H 是连通的（化归到 G 概单且 \(\Phi\) 循环的情形，并应用定理 1（节 3）以及第 VIII 章，§5，习题 13）。
\item
  假定 \(\mathbf{G}\) 是概单的。证明
\end{enumerate}

若 G 是 \(\mathrm{A}_{2l}\) 型（\(l \geq 1\)）的且 \(\varPhi\) 的阶为 2，则 \(\mathrm{H}_0\) 是 \(\mathrm{B}_l\) 型的；

若 G 是 \(A_{2l-1}\) 型（\(l \geq 2\)）的且 \(\Phi\) 的阶为 2，则 \(H_{0}\) 是 \(C_{l}\) 型的；

若 G 是 \(D_{l}\) 型 \((l \geq 4)\) 的且 \(\Phi\) 的阶为 2，则 \(H_{0}\) 是 \(B_{l-1}\) 型的；

若 G 是 \(D_{4}\) 型的且 \(\Phi\) 的阶为 3 或 6，则 \(H_{0}\) 是 \(G_{2}\) 型的；

若 \(\mathrm{G}\) 是 \(\mathrm{E}_6\) 型的且 \(\varPhi\) 的阶为 2，则 \(\mathrm{H}_0\) 是 \(\mathrm{F}_4\) 型的。

在每种情形中，确定群 \(\pi_i(\mathrm{H})\)，\(i = 0,1\)（参见~习题 3）。

\begin{enumerate}
\def\labelenumi{\alph{enumi})}
\setcounter{enumi}{2}
\tightlist
\item
  证明子群 \(\mathrm{H}_0\) 是主的（§4，习题 18）。
\end{enumerate}

\(d\) ) 证明 \(\mathrm{B}_3\) 或 \(\mathrm{A}_6\) 型的半单群包含一个 \(\mathrm{G}_2\) 型的主子群（利用 \(b\)) 以及 §4 的习题 20）。

\begin{enumerate}
\def\labelenumi{\arabic{enumi})}
\setcounter{enumi}{4}
\tightlist
\item
  假定群 G 是概单的；设 H 是 G 的一个主连通闭子群（§4，习题 18）。用 \(\Phi\) 表示 G 中使 H 固定的自同构 u 组成的群，用 F 表示 G 中被 \(\Phi\) 固定的元素组成的子群。
\end{enumerate}

\begin{enumerate}
\def\labelenumi{\alph{enumi})}
\item
  证明存在 \(\mathrm{G}\) 的一个在 \(\varPhi\) 下稳定的标架。
\item
  证明我们必定处于下列情形之一：
\end{enumerate}

\begin{enumerate}
\def\labelenumi{(\roman{enumi})}
\item
  \(H = F_{0};\)
\item
  G 是 \(B_{3}\) 型的，H 是 \(G_{2}\) 型的，且 \(\Phi\) 化为单位元；
\item
  G 是 \(A_{6}\) 型的，H 是 \(G_{2}\) 型的，且 \(\Phi\) 的阶为 2。
\end{enumerate}

（利用习题 4 以及 §4 的习题 20。）

¶ 6) 假定 G 是单连通的。设 p 是一个素数，g 是 C(G) 中满足 \(g^{p} = e\) 的一个元素。

\begin{enumerate}
\def\labelenumi{\alph{enumi})}
\tightlist
\item
  证明存在元素 \(u \in T\) 与 \(w \in W\)，使得
\end{enumerate}

\begin{enumerate}
\def\labelenumi{(\roman{enumi})}
\item
  \(w(u)u^{-1} = g;\)
\item
  \(w^{p} = 1;\)
\item
  \(u^{p} = e\) 当 \(p \neq 2\)，\(u^{p} = e\) 或 g 当 p = 2。
\end{enumerate}

（设 A 是 t 中的一个单房；对 \(i = 0, 1, \ldots, p - 1\)，设 \(x_i \in \overline{A}\) 满足 \(\exp x_i = g^i\)。取 u 为元素 \(\exp x\)，其中 x 是 \(\overline{A}\) 的以诸 \(x_i\) 为顶点的那个面的重心，并取 w 为 W 中满足 \(w(\mathrm{A}) = \mathrm{A} - x\) 的元素。）

\begin{enumerate}
\def\labelenumi{\alph{enumi})}
\setcounter{enumi}{1}
\tightlist
\item
  证明存在 \(u \in \mathrm{T}\) 与 \(v \in \mathrm{N}_{\mathrm{G}}(\mathrm{T})\)，使得
\end{enumerate}

\begin{enumerate}
\def\labelenumi{(\roman{enumi})}
\item
  \(vuv^{-1}u^{-1}=g;\)
\item
  \(v^{p} = e\) 当 \(p \neq 2\)，\(v^{p} = e\) 或 g 当 p = 2；
\item
  \(u^{p} = e\) 当 \(p \neq 2\)，\(u^{p} = e\) 或 \(g\) 当 \(p = 2\)。
\end{enumerate}

（利用 \(a\)) 中的构造，按 §4 习题 12 的方式把 \(w\) 提升为 \(n_w\)；取 \(v = n_w^{p+1}\) 当 \(p \neq 2\)，取 \(v = n_w\) 当 \(p = 2\)。）

\begin{enumerate}
\def\labelenumi{\arabic{enumi})}
\setcounter{enumi}{6}
\tightlist
\item
  设 p 是一个素数。证明下列条件等价：(i) p 不整除 \(\pi_{1}(G)\) 的挠子群的阶；
\end{enumerate}

\begin{enumerate}
\def\labelenumi{(\roman{enumi})}
\setcounter{enumi}{1}
\tightlist
\item
  对 G 的每个阶为 p 的元素 g，g 在 G 中的中心化子都是连通的；(iii) G 的每个同构于 \((\mathbf{Z}/p\mathbf{Z})^{2}\) 的子群都包含于一个极大环面。
\end{enumerate}

（为证明 (i) \(\Longrightarrow\) (ii)，利用习题 3 \(c\))；为证明 (iii) \(\Longrightarrow\) (i)，利用习题 6。）

\begin{enumerate}
\def\labelenumi{\arabic{enumi})}
\setcounter{enumi}{7}
\tightlist
\item
  取 \(\mathrm{G} = \mathbf{SO}(8, \mathbf{R}) / \{\pm I_8\}\)；用 \((\alpha_i)_{1 \leq i \leq 4}\) 表示 \(\mathrm{R}(\mathrm{G}, \mathrm{T})\) 的一个基，使得 \(\alpha_1, \alpha_3\) 与 \(\alpha_4\) 都不与 \(\alpha_2\) 正交（第 VI 章，图版 IV）。设 A 是 T 中由满足下列条件的 \(t \in T\) 组成的子群
\end{enumerate}

\[
\alpha_ {1} (t) = \alpha_ {3} (t) = \alpha_ {4} (t), \alpha_ {1} (t) ^ {2} = \alpha_ {2} (t) ^ {2} = 1.
\]

证明 A 同构于 \((\mathbf{Z}/2\mathbf{Z})^{2}\)，且它在 G 中的中心化子是一个有限的非交换群。

\begin{enumerate}
\def\labelenumi{\arabic{enumi})}
\setcounter{enumi}{8}
\tightlist
\item
  设 R 是一个不可约根系，\(R^{\vee}\) 是逆根系，B 是 R 的一个基，\(\alpha\) 是 R 的（相对于 B 的）最高根。置
\end{enumerate}

\[
\alpha = \sum_ {\beta \in \mathrm{B}} n _ {\beta} \beta \text { and } \alpha^ {\vee} = \sum_ {\beta \in \mathrm{B}} n _ {\beta} ^ {\vee} \beta^ {\vee};
\]

设 \(\nu (\mathrm{R}) = \sup_{\beta \in \mathrm{B}}n_{\beta}^{\vee}\)

\begin{enumerate}
\def\labelenumi{\alph{enumi})}
\item
  证明 N 中的区间 \(\mathbf{1},\nu(\mathrm{R})\mathbf{I}\) 是 1 与诸 \(n_{\beta}^{\vee}\)（\(\beta\in B\)）的并（置 \(\alpha_{0}=\alpha\)；设 \((\alpha_{1},\ldots,\alpha_{q})\) 是由 B 的互不相同的元素组成的一个序列，使得 \(\alpha_{i}\) 不与 \(\alpha_{i-1}\) 正交（对 \(i=1,\ldots,q\)），使得 \(n_{\alpha_{q}}^{\vee}=\nu(\mathrm{R})\)，且 q 对这些性质而言是极大的；证明 \(n_{\alpha_{i}}^{\vee}=i+1\)（对 \(i=1,\ldots,q\)））。
\item
  设 \(p\) 是一个素数。证明下列三条性质等价：
\end{enumerate}

\begin{enumerate}
\def\labelenumi{(\roman{enumi})}
\item
  \(p \leq \nu(\mathrm{R});\)
\item
  存在 \(\beta \in \mathrm{B}\) 使 \(p = n_{\beta}^{\vee}\)；
\item
  存在 \(\beta \in \mathrm{B}\) 使 \(p\) 整除 \(n_{\beta}^{\vee}\)。
\end{enumerate}

此时称 p 为 R 的一个挠素数。

\begin{enumerate}
\def\labelenumi{\alph{enumi})}
\setcounter{enumi}{2}
\tightlist
\item
  对不可约根系的每种型，给出 \(\nu(\mathrm{R})\) 的值以及挠素数。证明诸 \(n_{\beta}\) 组成的集合与诸 \(n_{\beta}^{\vee}\) 组成的集合除 \(G_{2}\) 型外是一致的。
\end{enumerate}

\(d\) ) 设 \(\mathbf{R}'\) 是 \(\mathbf{R}\) 的一个作为根系而言不可约的对称闭子集。证明 \(\nu (\mathrm{R}^{\prime})\leq \nu (\mathrm{R})\)

\begin{enumerate}
\def\labelenumi{\arabic{enumi})}
\setcounter{enumi}{9}
\tightlist
\item
  设 \(\mathbf{R}\) 是一个根系，\(p\) 是一个素数。证明下列性质等价：
\end{enumerate}

\begin{enumerate}
\def\labelenumi{(\roman{enumi})}
\item
  p 是 R 的某个不可约分支的挠素数（习题 9 b)）；\\
\item
  存在 R 的一个对称闭子集 \(R_{1}\)，它异于 R 且对这些性质而言是极大的，使得 \((\mathrm{Q}(\mathrm{R}^{\vee}):\mathrm{Q}(\mathrm{R}_{1}^{\vee}))=p(\mathrm{Q}(\mathrm{R}^{\vee}))\) 表示由 R 的逆根生成的 Z-模，而 \(\mathrm{Q}(\mathrm{R}_{1}^{\vee})\) 表示由诸 \(\alpha^{\vee}\)（\(\alpha\in R_{1}\)）生成的子模）；
\item
  存在 R 的一个对称闭子集 \(R_{1}\)，使 \(\mathrm{Q}(\mathrm{R}^{\vee})/\mathrm{Q}(\mathrm{R}_{1}^{\vee})\) 的 p-挠子模非零。
\end{enumerate}

（为证明 (i) \(\Longrightarrow\) (ii)，利用第 VI 章，§4 的习题 4；为证明 (iii) \(\Longrightarrow\) (i)，利用习题 9 d)。）

此时称 \(p\) 为 \(\mathbb{R}\) 的一个挠素数。

\begin{enumerate}
\def\labelenumi{\arabic{enumi})}
\setcounter{enumi}{10}
\tightlist
\item
  设 p 是一个素数。若存在 G 的一个具有极大秩的连通闭子群 H，使 \(\pi_{1}(H)\) 的 p-挠子群非零，则称 p 为 G 的一个挠素数。置 \(R = R(G, T)\)。
\end{enumerate}

\begin{enumerate}
\def\labelenumi{\alph{enumi})}
\item
  G 的挠素数是 R 的挠素数（习题 10）以及整除 \(\pi_1(\mathrm{G})\) 的挠群的阶的素数。
\item
  证明 G 的每个挠素数都整除 \(w/l!\)，其中 w 是 W 的阶，l 是半单群 D(G) 的秩（利用第 VI 章，§2，节 4，命题 7）。
\item
  设 \(t \in T\) 与 \(n \in N\) 满足 \(t^{n} \in \mathrm{C}(\mathrm{G})\)。设 \(R_{1}\) 是由 \(\alpha \in R\)（满足 \(\alpha(t) = 1\)）组成的集合，m 是 \(\mathrm{Q}(\mathrm{R}^{\vee}) / \mathrm{Q}(\mathrm{R}_{1}^{\vee})\) 的挠群的阶（参见~习题 10）。证明 m 的每个素因子都整除 n（化归到 G 单连通且概单的情形）。
\item
  假定 G 是单连通的。设 \(g \in G\) 与 \(n \in N\) 满足 \(g^{n}\) 属于 \(\mathrm{C}(\mathrm{G})\)，且 G 的任何挠素数都不整除 n。~证明 g 的中心化子的换位子群是单连通的。
\item
  设 g 是 G 的一个元素，p 是一个素数，满足 \(g^{p} = e\) 且 p 不是 G 的挠素数。证明中心化子 \(Z(g)\) 是连通的，且 p 不是 \(Z(g)\) 的挠素数。
\item
  假定 G 是单连通的，并设 p 是 G 的一个挠素数。证明存在 G 的一个阶为 p 的元素 g，使 \(\pi_{1}(\mathrm{Z}(g))\) 是 p 阶循环群（利用习题 10 以及第 VI 章，§4 的习题 4）。
\end{enumerate}

\begin{enumerate}
\def\labelenumi{\arabic{enumi})}
\setcounter{enumi}{11}
\tightlist
\item
  设 \(p\) 是一个素数。证明下列条件等价：
\end{enumerate}

\begin{enumerate}
\def\labelenumi{(\roman{enumi})}
\item
  \(p\) 不是 \(\mathbf{G}\) 的挠素数；
\item
  对每个子群 \(\mathbf{F}\)（属于 \(\mathbf{G}\)，同构于 \((\mathbf{Z} / p\mathbf{Z})^n\)（\(n\in \mathbf{N}\)）），\(\mathbf{F}\) 在 \(\mathbf{G}\) 中的中心化子都是连通的；
\end{enumerate}

(ii') 对 G 的每个同构于 \((\mathbf{Z}/p\mathbf{Z})^{2}\) 的子群 F，F 在 G 中的中心化子都是连通的；

\begin{enumerate}
\def\labelenumi{(\roman{enumi})}
\setcounter{enumi}{2}
\tightlist
\item
  G 的每个同构于 \((\mathbf{Z}/p\mathbf{Z})^{n}\)（n 为某个整数）的子群都包含于一个极大环面；
\end{enumerate}

(iii') G 的每个同构于 \((\mathbf{Z}/p\mathbf{Z})^{3}\) 的子群都包含于一个极大环面。

（为证明 (i) \(\Longrightarrow\) (ii)，利用习题 11 e)；为证明 (iii) \(\Longrightarrow\) (i)，利用习题 11 f) 以及习题 7。）

\exercisesection{\S~6}

\begin{enumerate}
\def\labelenumi{\arabic{enumi})}
\tightlist
\item
  设 R 是实向量空间 V 中的一个既约根系。给空间 V 赋以一个在 W(R) 下不变的内积，并给空间 S(V) 赋以相应的内积（《拓扑向量空间》第 V 章，§3，节 3，公式 (13) 与 (14)）；于是，对 V 中的 \(x_{1},\ldots,x_{n},y_{1},\ldots,y_{n}\)，
\end{enumerate}

\[
(x _ {1} \dots x _ {n} | y _ {1} \dots y _ {n}) = \sum_ {\sigma \in \mathfrak {S} _ {n}} (x _ {1} | y _ {\sigma (1)}) \dots (x _ {n} | y _ {\sigma (n)}).
\]

选取 R 的一个房 C；置 \(N = \text{Card}(R_{+})\)，\(\rho = \frac{1}{2} \sum_{\alpha \in R_{+}} \alpha\) 与 \(w(R) = \text{Card}(W(R))\)。设 P 是元素 \(\prod_{\alpha \in R_{+}} \alpha\)（\(\mathbf{S}^{\mathrm{N}}(\mathrm{V})\) 中的）。

\(a)\) 证明 \(\mathrm{P} = \frac{1}{\mathrm{N!}}\sum_{w\in \mathrm{W(R)}}\varepsilon (w)(w\rho)^{\mathrm{N}}\)（参见~第 VI 章，§3，节 3，命题 2）。

\begin{enumerate}
\def\labelenumi{\alph{enumi})}
\setcounter{enumi}{1}
\item
  推出等式 \((\mathrm{P}|\mathrm{P}) = w(\mathrm{R})\prod_{\alpha \in \mathrm{R}_{+}}(\rho |\alpha)\)。
\item
  证明 \((\mathrm{P}|\mathrm{P}) = 2^{-\mathrm{N}}w(\mathrm{R})\prod_{i = 1}^{l}m_i!\prod_{\alpha \in \mathrm{R}_+}(\alpha |\alpha) = 2^{-\mathrm{N}}\prod_{i = 1}^{l}(m_i + 1)!\prod_{\alpha \in \mathrm{R}_+}(\alpha |\alpha)\)，其中 \(m_{1},\ldots ,m_{l}\) 是 \(\mathrm{W(R)}\) 的指数（利用第 VIII 章，§9 的习题 3）。
\end{enumerate}

¶ 2) 设 V 是 R 上的一个有限维实希尔伯特空间。给 V 上的多项式空间 \(\mathbf{S}(\mathrm{V}^{*})\) 赋以《拓扑向量空间》第 V 章，§3，节 3，公式 (13) 与 (14) 所定义的内积（参见~习题 1）。用 \(\gamma\) 表示 V 上的典则高斯测度（《积分》第 IX 章，§6，节 4 至 6）；于是，若 \((x_{i})_{1\leq i\leq n}\) 是 \(V^{*}\) 的一个标准正交基，则有 \(d\gamma(x)=(2\pi)^{-n/2}e^{-(x|x)/2}dx_{1}\ldots dx_{n}\)。

\begin{enumerate}
\def\labelenumi{\alph{enumi})}
\tightlist
\item
  设 \(q\) 是 \(\mathbf{S}^2(\mathrm{V})\) 中定义 \(\mathrm{V}^*\) 上内积的元素，并设 \(\Delta: \mathbf{S}(\mathrm{V}^*) \to \mathbf{S}(\mathrm{V}^*)\) 是由与 \(q\) 作内积所给出的算子（《代数》第 III 章，§11，节 9）。证明对每个标准正交基 \((x_i)_{1 \leq i \leq n}\)（\(\mathrm{V}^*\) 的）以及每个多项式 \(\mathrm{P} \in \mathbf{S}(\mathrm{V}^*)\)，有 \(\Delta(\mathrm{P}) = \frac{1}{2} \sum_{i=1}^{n} \frac{\partial^2\mathrm{P}}{\partial x_i^2}\)。b) 对 \(\mathrm{P} \in \mathbf{S}(\mathrm{V}^*)\)，置 \(\mathrm{P}^* = \mathrm{P} * \gamma\)，使得 \(\mathrm{P}^*(x) = \int_{\mathrm{V}} \mathrm{P}(x - y) d\gamma(y)\)（对 \(x \in \mathrm{V}\)）。证明等式 \(\mathrm{P}^* = \sum_{n=1}^{\infty} \frac{\Delta^n(\mathrm{P})}{n!} = e^{\Delta}(\mathrm{P})\)。
\end{enumerate}

（化归为对函数 \(x \mapsto e^{i(x|u)}\)（\(u \in \mathrm{V}\)）证明类似的公式。）

\begin{enumerate}
\def\labelenumi{\alph{enumi})}
\setcounter{enumi}{2}
\tightlist
\item
  证明公式 \(\int\overline{\mathrm{P}^{*}(ix)}\mathrm{Q}^{*}(ix)d\gamma(x)=(\mathrm{P}|\mathrm{Q})\)（对 \(\mathbf{S}(\mathrm{V}^{*})\) 中的 P 与 Q，后者等同于 \(\mathbf{S}_{\mathbf{C}}((\mathrm{V}\otimes\mathbf{C})^{*})\) 的一个子空间）。
\end{enumerate}

\(d)\) 若 \(\mathrm{P}\) 是 \(\mathrm{V}\) 上的一个齐次多项式且 \(\Delta (\mathrm{P}) = 0\)，则 \(\int_{\mathrm{V}}\mathrm{P}(x)^2 d\gamma (x) = (\mathrm{P}|\mathrm{P})\)。

\begin{enumerate}
\def\labelenumi{\alph{enumi})}
\setcounter{enumi}{4}
\tightlist
\item
  设 W 是 V 的一个由反射生成的有限自同构群；用 H 表示 W 中反射组成的集合。对 \(h \in H\)，设 \(e_{h} \in V\) 与 \(f_{h} \in V^{*}\) 满足 \(h(x) = x + f_{h}(x)e_{h}\)。证明多项式 \(P = \prod_{h \in H} f_{h}\) 满足 \(\Delta(P) = 0\)（利用第 V 章，§5，节 4，命题 5）。
\end{enumerate}

\begin{enumerate}
\def\labelenumi{\arabic{enumi})}
\setcounter{enumi}{2}
\tightlist
\item
  设 \(\mu\) 是加法群 \(\mathfrak{g}\) 上的一个 Haar 测度。
\end{enumerate}

\begin{enumerate}
\def\labelenumi{\alph{enumi})}
\item
  G 上存在唯一的具有下列性质的 Haar 测度 \(\mu_{G}\)：若 \(\omega_{G}\) 与 \(\omega_{g}\) 分别是 G 与 g 上次数为 n 的不变微分形式，满足 \(\omega_{\mathrm{G}}(e)=\omega_{\mathfrak{g}}(0)\) 且 \(|\omega_{g}|=\mu\)，则 \(|\omega_{G}|=\mu_{G}\)。映射 \(\mu\mapsto\mu_{G}\) 是从 g 上 Haar 测度的集合到 G 上相应的集合的一个双射。
\item
  选取一个在 g 下不变的内积 ( \textbar{} )，以及 t 上的一个 Haar 测度 \(\tau\)，使得 \(\mu\)（相应地 \(\tau\)）当 g（相应地 t）借助一个标准正交基与 \(R^{n}\)（相应地与 \(R^{r}\)）等同起来时对应于 Lebesgue 测度。证明公式
\end{enumerate}

\[
\int_ {t} \pi_ {\mathfrak {g}} (x) e ^ {- (x | x) / 2} d \tau (x) = (2 \pi) ^ {n / 2} \frac {\tau_ {\mathrm{T}} (\mathrm{T})}{\mu_ {\mathrm{G}} (\mathrm{G})} w (\mathrm{G}).
\]

\begin{enumerate}
\def\labelenumi{\alph{enumi})}
\setcounter{enumi}{2}
\tightlist
\item
  把 X(T) 与希尔伯特空间 \(t^{*}\) 的一个子集借助映射 \((2\pi i)^{-1}\delta\)（§4，节 2）等同起来，并置 \(\mathrm{P}(x)=\prod_{\alpha\in\mathrm{R}_{+}}\langle\alpha,x\rangle\)（对 \(x\in t\)）。用习题 1 与 2 的记号，证明
\end{enumerate}

\[
\frac {\tau_ {\mathrm{T}} (\mathrm{T})}{\mu_ {\mathrm{G}} (\mathrm{G})} = (2 \pi) ^ {\mathrm{N}} w (\mathrm{G}) ^ {- 1} (\mathrm{P} | \mathrm{P}) = \pi^ {\mathrm{N}} \prod_ {\alpha \in \mathrm{R} _ {+}} (\alpha | \alpha) \prod_ {i = 1} ^ {l} m _ {i}!
\]

\(d\) ) 用 §3 习题 9 的记号，设 \(\mathfrak{g}_{\mathbf{Z}}\) 是 \(\mathfrak{g}\) 中由 \((2\pi)^{-1}\Gamma (\mathrm{T})\) 与诸元素 \(u_{\alpha},v_{\alpha}\)（\(\alpha \in \mathbb{R}\)）生成的 Z-李子代数。证明等式

\[
\mu_ {\mathrm{G}} (\mathrm{G}) = \mu (\mathfrak {g} / \mathfrak {g} _ {\mathbf {Z}}) \frac {2 ^ {r} \pi^ {\mathrm{N} + r}}{\prod_ {i} m _ {i} !}.
\]

\begin{enumerate}
\def\labelenumi{\alph{enumi})}
\setcounter{enumi}{4}
\tightlist
\item
  假定 G 是单连通的。证明 l = r，且
\end{enumerate}

\[
\mu_ {\mathrm{G}} (\mathrm{G}) = 2 ^ {l / 2} \pi^ {- \mathrm{N}} f ^ {1 / 2} \prod_ {\alpha \in \mathrm{R} _ {+}} (\alpha | \alpha) ^ {- 1} \prod_ {i = 1} ^ {l} (\alpha_ {i} | \alpha_ {i}) ^ {- 1 / 2} \prod_ {i} (m _ {i}!) ^ {- 1},
\]

其中 \(\{\alpha_1, \ldots, \alpha_l\}\) 是 \(R\) 的一个基，\(f\) 是 \(R\) 的连通指数。

\(f\) ) 再假定 \(\mathbb{R}\) 不可约且其所有根都等长；取 \(\mathfrak{g}\) 上的内积为 Killing 型的相反数。证明 \(\mu_{\mathrm{G}}(\mathrm{G}) = (2\pi)^{\mathrm{N} + r}(2h)^{n / 2}f^{1 / 2}\prod_{i}(m_i!)^{-1}\)。

\begin{enumerate}
\def\labelenumi{\arabic{enumi})}
\setcounter{enumi}{3}
\tightlist
\item
  设 \(\mathrm{X}\) 是一个 \(\mathrm{C}^{\infty}\) 类微分流形。在本习题及以下诸习题中，我们将用 \(\mathrm{H}(\mathrm{X})\) 简单地表示分次 \(\mathbf{R}\)-空间 \(\mathrm{H}(\Omega (\mathrm{X}))\)。
\end{enumerate}

\begin{enumerate}
\def\labelenumi{\alph{enumi})}
\item
  证明 \(\Omega(\mathrm{X})\) 是一个结合且反交换的分次微分代数（《代数学》（Algèbre）第 X 章第~183 页，习题 18）。由此推出 \(\mathrm{H}(\mathrm{X})\) 具有一个自然的结合且反交换的分次代数结构。
\item
  若 X 连通且维数为 p，则 \(\mathrm{H}^{i}(\mathrm{X}) = 0\)（对 i \textgreater{} p），且 \(\dim_{\mathbf{R}} \mathrm{H}^{0}(\mathrm{X}) = 1\)。此外，若 X 是紧的，则空间 \(\mathrm{H}^{p}(\mathrm{X})\) 的维数为 1（利用《微分流形与解析流形，结果》，11.2.4）。
\item
  设 Y 是另一个 \(C^{\infty}\) 类流形，\(f: X \to Y\) 是一个 \(C^{\infty}\) 类映射。映射 \(f^{*}: \Omega(Y) \to \Omega(X)\) 是一个复形态射（《微分流形与解析流形，结果》，8.3.5）；证明 \(\mathrm{H}(f^{*}) : \mathrm{H}(Y) \to \mathrm{H}(X)\) 是一个代数态射。
\item
  假定给定了 G 在 X 上的一个 \(C^{\infty}\) 类作用律。证明由不变形式组成的子复形 \(\Omega(\mathrm{X})^{\mathrm{G}}\) 是 \(\Omega(\mathrm{X})\) 的一个子代数。由此推出定理 2 中定义的映射 \(\mathrm{H}(i):\mathrm{H}(\Omega(\mathrm{X})^{\mathrm{G}})\to\mathrm{H}(\mathrm{X})\) 是分次代数的一个同构。
\item
  证明 \((\operatorname{Alt}(\mathfrak{g}))^{\mathrm{G}}\) 是 \(\operatorname{Alt}(\mathfrak{g})\) 的一个子代数，且分次代数 \(\operatorname{H}(\operatorname{G})\) 与之同构。
\end{enumerate}

\begin{enumerate}
\def\labelenumi{\arabic{enumi})}
\setcounter{enumi}{4}
\tightlist
\item
  用 H(G) 表示分次 R-代数 H(Ω(G))（参见~习题 4）。
\end{enumerate}

\(a\) ) 空间 \(\mathrm{H}^p (\mathbf{G})\) 当 \(p > n\) 时为零，当 \(p = n\) 与 \(p = 0\) 时维数为 1。空间 \(\mathrm{H}^1 (\mathrm{G})\) 可典则地与 \(\mathfrak{c}^*\) 等同，其中 \(\mathfrak{c} = \mathrm{L}(\mathrm{C}(\mathrm{G}))\)。

\begin{enumerate}
\def\labelenumi{\alph{enumi})}
\setcounter{enumi}{1}
\item
  以下假定 G 是半单的。证明 \(\mathrm{H}^{1}(\mathrm{G})=\mathrm{H}^{2}(\mathrm{G})=0\)（参见~第 I 章，§6，习题 1）。
\item
  用 \(\mathrm{B}(\mathfrak{g})\) 表示 \(\mathfrak{g}\) 上 G-不变的对称双线性型组成的空间；对 \(b \in \mathrm{B}(\mathfrak{g})\) 与 \(x, y, z\)（属于 \(\mathfrak{g}\)），置 \(\tilde{b}(x, y, z) = b([x, y], z)\)。证明映射 \(b \mapsto \tilde{b}\) 定义了从 \(\mathrm{B}(\mathfrak{g})\) 到 \(\mathrm{H}^{3}(\mathrm{G})\) 的一个同构（设 \(\omega\in\mathrm{H}^{3}(\mathrm{G})\)；证明对所有 \(x\in g\)，存在 g 上唯一的线性形式 \(f(x)\) 使 \(df(x)=i(x)\omega\)，并考察形式 \((x,y)\mapsto-\langle y,f(x)\rangle\)）。
\end{enumerate}

\(d\)) 证明 \(\mathbf{R}\)-向量空间 \(\mathrm{H}^3 (\mathrm{G})\) 的维数等于 \(\mathfrak{g}\) 的单理想的个数。

\begin{enumerate}
\def\labelenumi{\arabic{enumi})}
\setcounter{enumi}{5}
\tightlist
\item
  置 \(b_{i}(\mathrm{G}) = \dim_{\mathbf{R}} \mathrm{H}^{i}(\mathrm{G})\)（对 \(i \geq 0\)），且若 \(\mathrm{X}\) 表示一个不定元，置 \(\mathrm{P}_{\mathrm{G}}(\mathrm{X}) = \sum_{i \geq 0} b_{i}(\mathrm{G}) \mathrm{X}^{i}\)。
\end{enumerate}

\begin{enumerate}
\def\labelenumi{\alph{enumi})}
\item
  证明 \(b_{i}(\mathrm{G}) = \int_{\mathrm{G}} \mathrm{Tr} \wedge^{i} (\mathrm{Ad} g) dg\) 与 \(\mathrm{P}_{\mathrm{G}}(\mathrm{X}) = \int_{\mathrm{G}} \det(1 + \mathrm{X}. \mathrm{Ad} g) dg\)（参见~附录 II，引理 1）。
\item
  推出等式 \(\sum_{i} b_{i}(\mathrm{G}) = 2^{r}\)，以及 \(\sum_{i} (-1)^{i} b_{i}(\mathrm{G}) = 0\)（当 \(\dim \mathrm{G} > 0\) 时）（利用 H. Weyl 公式，或 §2 的习题 8）。
\item
  取 \(\mathrm{G} = \mathbf{U}(n, \mathbf{C})\)。证明 \(\mathrm{P}_{\mathrm{G}}(\mathrm{X})\) 是 \((\mathrm{X}_1 \ldots \mathrm{X}_n)^{2n-2}\) 在多项式 \(\frac{1}{n!}(1 + \mathrm{X})^n \prod_{\substack{1 \leq i, j \leq n \\ i \neq j}} (\mathrm{XX}_i + \mathrm{X}_j)(\mathrm{X}_i - \mathrm{X}_j)\)（系数在 \(\mathbf{Z}[\mathrm{X}]\) 中）中的系数。
\end{enumerate}

\begin{enumerate}
\def\labelenumi{\arabic{enumi})}
\setcounter{enumi}{6}
\tightlist
\item
  设 \(\mathbf{K}\) 是一个连通紧李群。
\end{enumerate}

\begin{enumerate}
\def\labelenumi{\alph{enumi})}
\item
  设 \(f: K \to G\) 是一个核为有限的满同态。证明同态 \(\mathrm{H}(f^{*}) : \mathrm{H}(\mathrm{G}) \to \mathrm{H}(\mathrm{K})\) 是一个同构。
\item
  证明代数 \(\mathrm{H}(\mathrm{G}\times\mathrm{K})\) 可典则地与斜张量积 \(\mathrm{H}(\mathrm{G})^{g}\otimes\mathrm{H}(\mathrm{K})\) 等同。
\item
  由 a) 与 b) 推出，作为代数，\(\mathrm{H}(\mathrm{G})\) 同构于 \(\mathrm{H}(\mathrm{C}(\mathrm{G})_{0})^{g}\otimes\mathrm{H}(\mathrm{D}(\mathrm{G}))\)。证明代数 \(\mathrm{H}(\mathrm{C}(\mathrm{G})_{0})\) 同构于 \(\wedge(\mathfrak{c}^{*})\)，其中 \(\mathfrak{c}=\mathrm{L}(\mathrm{C}(\mathrm{G}))\)。
\end{enumerate}

¶ 8) 设 k 是一个特征为零的域，E 是 k 上的一个分次左双代数（《代数》第 III 章，§11，节 4，定义 3）。假定 E 是反交换且反余交换的，且当 m 充分大时 \(E_{m} = 0\)。用 P 表示 E 的本原元组成的集合（参见~第 II 章，§1）。

\begin{enumerate}
\def\labelenumi{\alph{enumi})}
\item
  证明 P 的每个齐次元素都是奇次数的（对大的 m 展开 \(c(x^{m})\)）。由此推出存在分次左双代数的一个典则态射 \(\varphi : \Lambda(\mathrm{P}) \to \mathrm{E}\)（\(\Lambda(\mathrm{P})\) 的分次左双代数结构即《代数》第 III 章，§11，习题 6 中所定义的那一个）。
\item
  证明 \(\varphi\) 是一个同构（把第 II 章，§1，节 6 定理 1 的证明加以改造）。
\end{enumerate}

\begin{enumerate}
\def\labelenumi{\arabic{enumi})}
\setcounter{enumi}{8}
\tightlist
\item
  用 \(m: \mathrm{G} \times \mathrm{G} \to \mathrm{G}\) 表示满足 \(m(g, h) = gh\)（对 \(g, h\)，属于 \(\mathrm{G}\)）的映射。把 \(\mathrm{H}(\mathrm{G} \times \mathrm{G})\) 与 \(\mathrm{H}(\mathrm{G})^g \otimes \mathrm{H}(\mathrm{G})\) 等同起来（习题 7 b)），于是 \(m^*\) 定义了一个代数同态 \(c: \mathrm{H}(\mathrm{G}) \to \mathrm{H}(\mathrm{G})^g \otimes \mathrm{H}(\mathrm{G})\)。
\end{enumerate}

\begin{enumerate}
\def\labelenumi{\alph{enumi})}
\item
  证明 \((\mathrm{H}(\mathrm{G}), c)\) 是一个反交换且反余交换的分次左双代数（注意映射 \(g \mapsto g^{-1}\) 在 \(\mathrm{H}^{p}(\mathrm{G})\) 上诱导乘以 \((-1)^{p}\)）。
\item
  设 \(\mathrm{P}(\mathrm{G})\) 是 \(\mathrm{H}(\mathrm{G})\) 中由本原元素组成的分次子空间；由习题 8 推出存在分次双代数的一个同构 \(\Lambda(\mathrm{P}(\mathrm{G})) \to \mathrm{H}(\mathrm{G})\)。
\item
  证明 \(\dim_{\mathbf{R}}\mathrm{P}(\mathrm{G})=r\)（利用习题 6 b)）。
\end{enumerate}

\(d\) ) 推出多项式 \(\sum_{i\geq 0}b_i(\mathrm{G})\mathrm{X}^i\) 具有如下形状

\[
(1 + \mathrm{X}) ^ {c} \prod_ {i = 1} ^ {l} (1 + \mathrm{X} ^ {2 k _ {i} + 1}),
\]

其中 c 是 \(\mathrm{C}(\mathrm{G})\) 的维数，l 是 \(\mathrm{D}(\mathrm{G})\) 的秩，诸 \(k_{i}\) 是 \(\geq 1\) 的整数；我们有 \(k_{1} + \cdots + k_{l} = \frac{1}{2}\mathrm{Card}\mathrm{R}(\mathrm{G},\mathrm{T})\)\(^{12}\)。

\footnotetext[12]{整数 \(k_{i}\) 事实上是 \(\mathrm{R}(\mathrm{G},\mathrm{T})\) 的指数。参见 J. LERAY, Sur l'homologie des groupes de Lie, des espaces homogènes et des espaces fibrés principaux, Colloque de Topologie de Bruxelles (1950), pp. 101-115.}

\begin{enumerate}
\def\labelenumi{\arabic{enumi})}
\setcounter{enumi}{9}
\tightlist
\item
  设 \(\mathrm{H}\) 是 \(\mathrm{G}\) 的一个闭子群；用 \(dh\) 表示 \(\mathrm{H}\) 上总质量为 1 的 Haar 测度。置 \(\chi(\mathrm{G} / \mathrm{H}) = \sum_{i > 0} (-1)^i \dim_{\mathbf{R}} \mathrm{H}^i (\mathrm{G} / \mathrm{H})\)。
\end{enumerate}

\begin{enumerate}
\def\labelenumi{\alph{enumi})}
\item
  证明等式 \(\chi(\mathrm{G}/\mathrm{H})=\overline{\int_{\mathrm{H}}}\det(1-\mathrm{Ad}_{\mathfrak{g}/\mathfrak{h}}h)\,dh\)（利用《代数学》（Algèbre）第 X 章第~41 页，命题 11 以及附录 II，节 2，引理 1）。
\item
  设 \(\pi_{0}(\mathrm{H})\) 是 H 的连通分支的个数。证明
\end{enumerate}

\(\chi (\mathrm{G / H}) = 0\)，若 \(\mathrm{H_0}\) 不是极大秩

\(\chi (\mathrm{G / H}) = w(\mathrm{G}) / w(\mathrm{H}_0)\pi_0(\mathrm{H})\)，若 \(\mathrm{H_0}\) 是极大秩。

\begin{enumerate}
\def\labelenumi{\arabic{enumi})}
\setcounter{enumi}{10}
\tightlist
\item
  设 \(u\) 是 \(G\) 的一个阶为 2 的自同构；用 \(K\) 表示 \(u\) 的不动点集合的单位元连通分支，\(\mathfrak{k}\) 表示它的李代数，\(X\) 表示对称空间 \(G / K\)（§1，习题 8）。
\end{enumerate}

\begin{enumerate}
\def\labelenumi{\alph{enumi})}
\item
  证明 X 上每个 G-不变微分形式 \(\omega\) 都满足 \(d\omega = 0\)（注意 \(u\) 在 \(\mathrm{Alt}^p (\mathfrak{g} / \mathfrak{k})\) 上诱导乘以 \((-1)^p)\)）。
\item
  由此推出，从分次代数 \(\mathrm{H}(\mathrm{G}/\mathrm{K})\) 到由 K-不变元素组成的 \(\operatorname{Alt}(\mathfrak{g}/\mathfrak{k})\) 的分次子代数有一个同构。
\item
  置 \(b_{i}(\mathrm{G / K}) = \dim_{\mathbf{R}}\mathrm{H}^{i}(\mathrm{G / K})\)（对 \(i\geq 0\)）；对 \(k\in \mathrm{K}\)，用 \(\mathrm{Ad}^{-}k\) 表示 \(\mathrm{Ad}k\) 在 \(\mathrm{L}(u)\) 的特征值为 \(-1\) 的特征子空间上的限制。设 \(dk\) 是 \(\mathbf{K}\) 上总质量为 1 的 Haar 测度。证明公式 \(b_{i}(\mathrm{G / K}) = \int_{\mathrm{K}}\mathrm{Tr}\wedge^{i}(\mathrm{Ad}^{-}k)dk\) 与 \(\sum_{i > 0}b_{i}(\mathrm{G / K})\mathrm{X}^{i} = \int_{\mathrm{K}}\det (1 + \mathrm{XAd}^{-}k)dk\)。
\item
  证明若此外 K 还是极大秩的，则代数 \(\mathrm{H}(\mathrm{G}/\mathrm{K})\) 的奇数部分为零（注意 \(u = \operatorname{Int} k\)，其中 \(k \in K\)）。
\item
  计算分次代数 \(\mathrm{H}(\mathbf{S}_{n})\)。由此推出 \(S_{n}\) 具有一个李群结构（与其流形结构相容）当且仅当 n 等于 1 或 3。
\end{enumerate}

\begin{enumerate}
\def\labelenumi{\arabic{enumi})}
\setcounter{enumi}{11}
\tightlist
\item
  设 H 是 G 的一个连通闭子群，h 是它的李代数。
\end{enumerate}

\begin{enumerate}
\def\labelenumi{\alph{enumi})}
\item
  设 \(\alpha\) 是 \(\mathfrak{h}^*\) 的一个在 H 下不变的元素；证明存在元素 \(\bar{\alpha}\)（属于 \(\mathfrak{g}^*\)，在 H 下不变，且它到 \(\mathfrak{h}\) 上的限制等于 \(\alpha\)）。证明元素 \(d\bar{\alpha} \in \mathrm{Alt}^2(\mathfrak{g})\) 被 \(i(\eta)\) 与 \(\theta(\eta)\)（对所有 \(\eta \in \mathfrak{h}\)）零化，且它在 \(\mathrm{H}^2(\mathrm{G}/\mathrm{H})\) 中的类不依赖于 \(\bar{\alpha}\) 的选取。这就定义了一个同态 \(\varphi: \mathrm{H}^1(\mathrm{H}) \to \mathrm{H}^2(\mathrm{G}/\mathrm{H})\)。
\item
  以下假定 \(\mathbf{G}\) 是半单的。证明 \(\varphi\) 是一个同构。
\item
  证明 \(\mathrm{H}\) 是半单的当且仅当 \(\mathrm{H}^2 (\mathrm{G / H}) = 0\)。
\end{enumerate}

\(d\) ) 不假定 \(\mathrm{H}\) 连通，定义一个同构 \((\mathfrak{h}^{*})^{\mathrm{H}}\to\) \(\mathrm{H}^2 (\mathrm{G / H})\)（把 \(b\) ) 应用于 \(\mathrm{H}_0\)）。

\begin{enumerate}
\def\labelenumi{\arabic{enumi})}
\setcounter{enumi}{12}
\tightlist
\item
  设 \(\mathrm{H}\) 是 \(\mathrm{G}\) 的一个闭子群；置 \(\mathrm{X} = \mathrm{G} / \mathrm{H}\)，\(n = \dim \mathrm{X}\)。证明下列条件等价：
\end{enumerate}

\begin{enumerate}
\def\labelenumi{(\roman{enumi})}
\item
  在 X 上存在一个 2-形式 \(\omega\)，使得 \(d\omega = 0\)，且使得交错形式 \(\omega_{x}\)（\(\mathrm{T}_x(\mathrm{X})\) 上的）对所有 \(x\in \mathrm{X}\) 都是分离的；
\item
  \(n\) 是偶数，且存在元素 \(\omega\)（属于 \(\mathrm{H}^2 (\mathrm{G / H})\)）使 \(\omega^{n / 2}\neq 0\)；(iii) \(\mathrm{H}\) 是 \(\mathbf{G}\) 中一个环面的中心化子；
\item
  H 是极大秩的，且在 X 上存在一个 G-不变的复结构；
\item
  在 X 上存在一个复结构 \(j\) 与一个 2-形式 \(\omega\)，使得 \(d\omega = 0\)，\(\omega_{x}(ju,jv) = \omega_{x}(u,v)\)，且 \(\omega_{x}(u,ju) > 0\)（对所有 X 中的 \(x\) 以及所有非零的 \(u,v\)（属于 \(\mathrm{T}_x(\mathrm{X})\)））（* 换言之，即 \(\mathrm{X}_{*}\) 上的一个 Kähler 结构*）；
\item
  在 X 上存在满足 (v) 中条件且在 G 下不变的复结构 j 与 2-形式 \(\omega\)（* 即 \(X_{*}\) 上的一个 G-不变 Kähler 结构*）。
\end{enumerate}

（为证明 (ii) \(\Longrightarrow\) (iii)，置 \(S = C(H)_{0}\) 与 \(Z = Z_{G}(S)\)；利用习题 12 d)，证明典则映射 \(H^{2}(G/Z) \to H^{2}(G/H)\) 是满射，并由此推出 Z = H。(iii) 与 (iv) 的等价性由 §4 的习题 8 得出。为证明 (iii) \(\Longrightarrow\) (vi)，在 g/L(H) 上构造一个 H-不变的分离正埃尔米特形式并考察其虚部；利用习题 11 a)。）

\exercisesection{\S~7}

\begin{enumerate}
\def\labelenumi{\arabic{enumi})}
\tightlist
\item
  取 G 为群 \(\mathbf{U}(n,\mathbf{C})\)，T 为由对角矩阵组成的子群；对 T 中的 \(t = \text{diag}(t_1, \ldots, t_n)\) 与 \(1 \leq i \leq n\)，置 \(\varepsilon_i(t) = t_i\)。用 \(\sigma\) 表示 G 在 \(C^n\) 上的恒等表示。
\end{enumerate}

\begin{enumerate}
\def\labelenumi{\alph{enumi})}
\item
  群 X(T) 有基 \(\varepsilon_{1},\ldots,\varepsilon_{n}\)；证明 \(\mathbf{Z}[\mathrm{X}(\mathrm{T})]^{\mathrm{W}}\) 的每个元素都具有形状 \(e^{k(\varepsilon_{1}+\cdots+\varepsilon_{n})}\mathrm{P}(e^{\varepsilon_{1}},\ldots,e^{\varepsilon_{n}})\)，其中 k 是一个相对整数，P 是 n 个变量的、系数为整数的对称多项式。
\item
  证明表示 \(\Lambda^r\sigma\)（\(1 \leq r \leq n\)）都是不可约的。
\item
  证明同态 \(u: \mathbf{Z}[X_1, X_2, \ldots, X_n][X_n^{-1}] \to \mathrm{R}(\mathrm{G})\)（满足 \(u(X_i) = [\Lambda^i \sigma]\)）是一个同构。
\end{enumerate}

\begin{enumerate}
\def\labelenumi{\arabic{enumi})}
\setcounter{enumi}{1}
\tightlist
\item
  设 \(\mathrm{G} = \mathbf{SO}(2l + 1, \mathbf{R})\)，其中 \(l \geq 1\)；我们使用 §4 习题 2 的记号。用 \(\sigma\) 表示 \(\mathrm{G}\) 在 \(\mathbf{C}^{2l + 1}\) 上由恒等表示通过标量扩张得到的表示。
\end{enumerate}

\begin{enumerate}
\def\labelenumi{\alph{enumi})}
\item
  Z-模 X(T) 以 \(\varpi_{1},\ldots,\varpi_{l-1},2\varpi_{l}\) 为基，同样也以 \(\varepsilon_{1},\ldots,\varepsilon_{l}\) 为基。
\item
  用 \(\eta_{i}\) 表示元素 \(e^{\varepsilon_{i}} + e^{-\varepsilon_{i}}\)（属于 Z{[}X(T){]}）。证明 \(\mathbf{Z}[X(T)]^{W}\) 的每个元素都可写成 \(\mathrm{P}(\eta_{1}, \ldots, \eta_{l})\)，其中 P 是 l 个变量的、系数为整数的对称多项式。
\item
  证明表示 \(\Lambda^r\sigma (r\leq 2l + 1)\) 都是不可约的。对 \(r\leq l\)，证明等式
\end{enumerate}

\[
\begin{array}{c} \operatorname{Ch} (\Lambda^ {r}   \sigma) = s _ {r} (\eta_ {1}, \ldots , \eta_ {l}) + (l - r + 2) s _ {r - 2} (\eta_ {1}, \ldots , \eta_ {l}) + \dots + \\ \qquad \qquad + \binom {l - r + 2 k} {k} s _ {r - 2 k} (\eta_ {1}, \ldots , \eta_ {l}) + \dots , \end{array}
\]

其中诸 \(s_{k}\) 是 l 个变量的初等对称多项式。

\(d\) ) 证明同态 \(u:\mathbf{Z}[\mathrm{X}_1,\dots ,\mathrm{X}_l]\to \mathrm{R}(\mathrm{G})\)（满足 \(u(\mathrm{X}_i) = [\Lambda^i\sigma ]\)）是一个同构。

\begin{enumerate}
\def\labelenumi{\arabic{enumi})}
\setcounter{enumi}{2}
\tightlist
\item
  设 \(\mathrm{G} = \mathbf{SO}(2l, \mathbf{R})\)，其中 \(l \geq 2\)；我们使用 §4 习题 2 的记号。用 \(\sigma\) 表示 G 在 \(C^{2l}\) 上由恒等表示通过标量扩张得到的表示。
\end{enumerate}

\begin{enumerate}
\def\labelenumi{\alph{enumi})}
\item
  \(\mathbf{Z}\)-模 \(\mathrm{X(T)}\) 有基 \(\varepsilon_1, \ldots, \varepsilon_l\)。
\item
  用 \(\eta_i\) 表示元素 \(e^{\varepsilon_i} + e^{-\varepsilon_i}\)（属于 \(\mathbf{Z}[\mathrm{X(T)}]\)），用 \(\delta\) 表示元素 \(\prod_{i=1}^{l}(e^{\varepsilon_i} - e^{-\varepsilon_i})\)。证明 \(\mathbf{Z}[\mathrm{X(T)}]^{\mathrm{W}}\) 的每个元素都可写成 \(\mathrm{P}(\eta_1, \ldots, \eta_l) + \mathrm{Q}(\eta_1, \ldots, \eta_l)\delta\)，其中 \(\mathrm{P}\) 与 \(\mathrm{Q}\) 是 \(l\) 个变量的、系数在 \(\frac{1}{2}\mathbf{Z}\) 中的对称多项式。
\item
  证明表示 \(\Lambda^{r}\sigma\) 当 \(r\neq l\) 时都是不可约的；表示 \(\Lambda^{l}\sigma\) 是两个子表示 \(\tau^{+}\) 与 \(\tau^{-}\) 的直和，其最高权分别为 \(2\varpi_{l}\) 与 \(2\varpi_{l-1}\)（见第 VIII 章，§13，习题 10）。
\end{enumerate}

\(d\) ) 证明当 \(r \leq l\) 时，元素 \(\operatorname{Ch}(\Lambda^r \sigma)\) 由习题 2 \(c\)) 中的公式给出，且

\[
\mathrm{Ch} (\tau^ {+}) = \frac {1}{2} (\delta + \mathrm{Ch} (\Lambda ^ {l} \sigma)) \qquad \mathrm{Ch} (\tau^ {-}) = \frac {1}{2} (- \delta + \mathrm{Ch} (\Lambda ^ {l} \sigma)).
\]

\begin{enumerate}
\def\labelenumi{\alph{enumi})}
\setcounter{enumi}{4}
\tightlist
\item
  证明同态 \(u : Z[X_{1}, \ldots, X_{l}; Y] \to R(G)\)（满足 \(u(X_{i}) = [\Lambda^{i} \sigma]\) 与 \(u(Y) = [\tau^{+}]\)）是满射，且其核由 \(Y^{2} - YX_{l} + A\) 生成，其中 \(A \in Z[X_{1}, \ldots, X_{l}]\)。
\end{enumerate}

\begin{enumerate}
\def\labelenumi{\arabic{enumi})}
\setcounter{enumi}{3}
\tightlist
\item
  设 E 是一个复向量空间，\(\mathbf{T}_{q}^{p}(\mathrm{E})\) 是 E 上 \((p,q)\) 型张量组成的空间（《代数》第 III 章，§5，节 6）。用 \(\mathrm{H}_{q}^{p}(\mathrm{E})\) 表示 \(\mathbf{T}_{q}^{p}(\mathrm{E})\) 中由对称张量（即属于 \(\mathbf{T}\mathbf{S}^{p}(\mathrm{E})\otimes\mathbf{T}\mathbf{S}^{q}(\mathrm{E}^{*})\) 在 \(\mathrm{T}_{q}^{p}(\mathrm{E})\) 中的像、且被收缩 \(c_{j}^{i}\)（对 \(i\in\{1,p\}\) 与 \(j\in\{p+1,p+q\}\)）零化的那些张量）组成的子空间（同上）。
\end{enumerate}

\begin{enumerate}
\def\labelenumi{\alph{enumi})}
\item
  证明 \(\mathrm{H}_{q}^{p}(\mathbf{C}^{n})\) 在 \(\mathbf{GL}(n,\mathbf{C})\) 下稳定，从而在 \(\mathbf{SU}(n,\mathbf{C})\) 下稳定；把 \(\mathbf{SU}(n,\mathbf{C})\) 的这个表示记作 \(\tau_{q}^{p}\)。
\item
  证明 \(\tau_q^p\) 是一个不可约表示，其最高权（用 §4 习题 1 的记号）是 \(p\varpi_1 + q\varpi_{n - 1}\)。
\item
  \(\mathbf{SU}(3,\mathbf{C})\) 的每个不可约表示都同构于诸表示 \(\tau_q^p\) 之一。
\end{enumerate}

\(d\) ) 设 \((x_{i})_{1\leq i\leq n}\) 是 \(\mathbf{C}^n\) 的典则基，\((y_{j})_{1\leq j\leq n}\) 是对偶基。证明 \(\mathrm{H}_q^p (\mathbf{C}^n)\) 可与由多项式 \(\mathrm{P}\in \mathbf{C}[x_1,\ldots ,x_n,y_1,\ldots ,y_n]\)（次数为 \(p\)，关于诸 \(x_{i}\)；次数为 \(q\)，关于诸 \(y_{i}\)；且满足 \(\sum_{i = 1}^{n}\frac{\partial^2\mathrm{P}}{\partial x_i\partial y_i} = 0\)）组成的空间等同。

\begin{enumerate}
\def\labelenumi{\arabic{enumi})}
\setcounter{enumi}{4}
\tightlist
\item
  设 \(k\) 是一个特征为零的交换域，\(V\) 是 \(k\) 上的一个向量空间，而 \(\varPsi\) 是 \(V\) 上一个分离的二次型。设 \(\varGamma \in \mathbf{S}^2(V^*)\)（相应地 \(\varGamma^* \in \mathbf{S}^2(V)\)）是伴随于 \(\varPsi\) 的元素（相应地 \(\varPsi\) 的逆型）。用 \(Q\) 表示 \(S(V)\) 上取与 \(\varGamma^*\) 的乘积所给出的自同态，用 \(\varDelta\) 表示 \(S(V)\) 上取与 \(\varGamma\) 的内积所给出的自同态，用 \(h\) 表示 \(S(V)\) 上在 \(S^r(V)\) 上约为乘以 \(-\frac{n}{2} - r\) 的自同态。\(a\) ) 若 \((x_i)_{1\leq i\leq n}\) 是 \(V\) 的一个标准正交基，则 \(Q(P) = \frac{1}{2}\left(\sum x_i^2\right).P\) 与 \(\varDelta(P) = \frac{1}{2}\sum \frac{\partial^2P}{\partial x_i^2}\)（对 \(P \in S(V)\)）。
\end{enumerate}

\begin{enumerate}
\def\labelenumi{\alph{enumi})}
\setcounter{enumi}{1}
\item
  证明公式 \([\Delta, \mathrm{Q}] = -h\)，\([h, \Delta] = 2\Delta\)，\([h, \mathrm{Q}] = -2\mathrm{Q}\)。
\item
  设 \(H_{r}\) 是 \(\mathbf{S}^{r}(\mathrm{V})\) 中被 \(\Delta\) 零化的子空间（「r 次齐次调和多项式」）。由 b) 推出存在一个在 \(\mathbf{O}(\varPsi)\) 下稳定的直和分解
\end{enumerate}

\[
\mathbf {S} ^ {r} (\mathrm{V}) = \mathrm{H} _ {r} \oplus \mathrm{QH} _ {r - 2} \oplus \mathrm{Q} ^ {2} \mathrm{H} _ {r - 4} \oplus \dots .
\]

\(d\) ) 证明表示 \(\mathrm{H}_r\) 是不可约的（参见~第 VIII 章，§13，节 3，(IV)）。

\begin{enumerate}
\def\labelenumi{\alph{enumi})}
\setcounter{enumi}{4}
\tightlist
\item
  取 \(k = \mathbf{C}\)，\(V = \mathbf{C}^n\) 带通常的二次型（\(n \geq 3\)）。我们于是得到不可约表示 \(\tau_r\)（\(\mathbf{SO}(n, \mathbf{R})\) 的）；用 §4 习题 2 的记号，证明 \(\tau_r\) 的最高权是 \(r\varpi_1\)。
\end{enumerate}

\(f\) ) 设 \(\Gamma\) 是由 Killing 型得到的 G 的 Casimir 元。证明公式 \(\Gamma_{\mathbf{S}(\mathrm{V})} = \frac{1}{2n - 4}\left(-4\mathrm{Q}\Delta +\left(\mathrm{H} + \frac{n}{2}\mathrm{I}\right)\left(\mathrm{H} + \left(2 - \frac{n}{2}\right)\mathrm{I}\right)\right)\)。计算 \(\tilde{\Gamma} (\tau_r)\) 并由此推出型 \(Q_{\Gamma}\) 的值（命题 4）。

\begin{enumerate}
\def\labelenumi{\arabic{enumi})}
\setcounter{enumi}{5}
\item
  假定 G 是概单的。证明 G 具有一个忠实不可约表示当且仅当它不同构于 \(\mathbf{Spin}(4k, \mathbf{R})\)（\(k \geq 2\)）。
\item
  设 \(\tau : \mathrm{G} \to \mathbf{SO}(n, \mathbf{R})\) 是 \(\mathrm{G}\) 的一个实酉表示；用 \(\varphi : \mathbf{Spin}(n, \mathbf{R}) \to \mathbf{SO}(n, \mathbf{R})\) 表示典则二重覆盖。于是称 \(\tau\) 为旋量的，若存在一个态射 \(\tilde{\tau} : \mathrm{G} \to \mathbf{Spin}(n, \mathbf{R})\) 使 \(\varphi \circ \tilde{\tau} = \tau\)。
\end{enumerate}

\begin{enumerate}
\def\labelenumi{\alph{enumi})}
\item
  设 \(\Sigma\) 是 \(\mathrm{P}(\mathrm{T},\tau)\) 的一个满足 \(\Sigma\cup(-\Sigma)=\mathrm{P}(\mathrm{T},\tau)\) 与 \(\Sigma\cap(-\Sigma)=\varnothing\) 的子集；用 \(\omega_{\Sigma}\) 表示 \(\Sigma\) 的诸元素之和。类 \(\varpi\)（\(\omega_{\Sigma}\) 在 \(\mathrm{X}(\mathrm{T})/2\mathrm{X}(\mathrm{T})\) 中的）不依赖于 \(\Sigma\) 的选取。证明 \(\tau\) 是旋量的当且仅当 \(\varpi=0\)。
\item
  证明 \(\rho \in \mathrm{X}(\mathrm{T})\) 当且仅当伴随表示是旋量的。
\end{enumerate}

\begin{enumerate}
\def\labelenumi{\arabic{enumi})}
\setcounter{enumi}{7}
\tightlist
\item
  设 G 是一个李代数为紧的、且只有有限个连通分支的李群。证明 G 在一个有限维向量空间上具有一个忠实线性表示（把 G 写成一个紧群 K 与一个向量群 N 的半直积；选取 K 在一个有限维向量空间 W 上的一个忠实表示，并把 G 表示为 \(W \oplus N\) 的仿射群的一个子群）。
\end{enumerate}

\exercisesection{\S~8}

\begin{enumerate}
\def\labelenumi{\arabic{enumi})}
\tightlist
\item
  设 \(\mathbf{G} = \mathbf{SU}(2, \mathbf{C})\)，\(\sigma\) 是 G 在 \(C^{2}\) 上的恒等表示。
\end{enumerate}

\begin{enumerate}
\def\labelenumi{\alph{enumi})}
\item
  G 的不可约表示是诸表示 \(\tau^{n} = S^{n}\sigma\)（\(n \geq 0\)）。
\item
  设 \(e_{1}, e_{2}\) 是 \(C^{2}\) 的典则基；证明 \(\tau^{n}\) 在基 \((e_{1}^{i}e_{2}^{n-i})_{0\leq i\leq n}\) 下的系数是满足如下条件的函数 \(\tau_{ij}^{n}\)：对 \(g=\begin{pmatrix}\alpha&\beta\\-\bar{\beta}&\bar{\alpha}\end{pmatrix}\in G\)，有 \(\tau_{ij}^{n}(g)=\frac{(-1)^{i}}{j!}\alpha^{i+j-n}\beta^{j-i}\mathrm{P}_{ij}^{n}(|\alpha|^{2})\)，其中 \(\mathrm{P}_{ij}^{n}(t)=\frac{d^{j}}{dt^{j}}\left[t^{n-i}(1-t)^{i}\right]\)（「Jacobi 多项式」）。
\item
  由 \(b)\) 推出，诸函数 \((n + 1)^{1/2}\left(\frac{j!(n - j)!}{i!(n - i)!}\right)^{1/2}\tau_{ij}^n\)（对整数 \(i,j,n\)，满足 \(0\leq i\leq n\)，\(0\leq j\leq n\)）构成 \(\mathrm{L}^2 (\mathrm{G})\) 的一个标准正交基。
\end{enumerate}

\(d)\) 对 \(g = \left( \begin{array}{cc}\alpha & \beta \\ -\bar{\beta} & \bar{\alpha} \end{array} \right)\in \mathrm{G}\)，置 \(\alpha = t^{1 / 2}e^{i(\varphi -\psi) / 2},\beta = (1 - t)^{1 / 2}e^{i(\varphi +\psi) / 2}\)，其中 \(0\leq t\leq 1,0\leq \varphi < 2\pi , - 2\pi \leq \psi < 2\pi\)。证明规范化的 Haar 测度 \(dg\)（\(\mathrm{G}\) 上的）等于 \((8\pi^2)^{-1}dtd\varphi d\psi\)。

\begin{enumerate}
\def\labelenumi{\alph{enumi})}
\setcounter{enumi}{4}
\tightlist
\item
  设 \(a, b\)（属于 \(\frac{1}{2}\mathbf{Z}\)）满足 \(a - b \in \mathbf{Z}\)。由 \(d\) ) 推出，多项式 \(\mathrm{P}_{n/2 - a, n/2 - b}^n(t)\)（\(n\) 为与 \(2a\) 同奇偶的整数，且 \(n \geq \max(2a, 2b)\)）构成 \(\mathrm{L}^2([0,1])\) 关于测度 \(t^{-a - b}(1 - t)^{a - b}dt\) 的一个正交基。
\end{enumerate}

\begin{enumerate}
\def\labelenumi{\arabic{enumi})}
\setcounter{enumi}{1}
\item
  设 \(f\) 是 \(G\) 上的一个 \(C^\infty\) 类复值函数。证明存在两个复值函数 \(g\) 与 \(\varphi\)（\(G\) 上的，\(C^\infty\) 类），使得 \(\varphi\) 是中心的且 \(f = g * \varphi\)。
\item
  设 u 是 G 的一个不可约表示，\(\lambda\) 是它的最高权。对 \(x \in g\)，用 \(\bar{x}\) 表示 \(\overline{C}\) 中在 \(\mathrm{Ad}(G)\) 下与 x 共轭的唯一元素。证明等式 \(\|u(x)\|_{\infty} = |\delta(\lambda)(\bar{x})|\)。
\item
  选取一个分离的正二次型 \(\mathbf{Q}\)（\(\mathfrak{g}\) 上的，在 \(\operatorname{Ad}(\mathbf{G})\) 下不变）。对 \(x \in \mathfrak{t}\)，置 \(\vartheta_0(x) = \sum_{u \in \Gamma(\mathbf{T})} e^{-\mathbf{Q}(x + u)}\)。
\end{enumerate}

\begin{enumerate}
\def\labelenumi{\alph{enumi})}
\item
  证明 \(\vartheta_0\) 是一个 \(C^\infty\) 类函数（在 \(\mathfrak{t}\) 上），且存在函数 \(\vartheta_1\)（\(C^\infty\) 类的，在 \(T\) 上）使 \(\vartheta_1 \circ \exp_T = \vartheta_0\)。
\item
  证明存在唯一的中心函数 \(\vartheta\)（\(\mathbf{G}\) 上的，\(\mathrm{C}^{\infty}\) 类），它在 \(\mathrm{T}\) 上的限制等于 \(\vartheta_{1}\)。对每个极大环面 \(\mathrm{S}\)（\(\mathrm{G}\) 的）与所有 \(x\in \mathrm{L}(\mathrm{S})\)，我们有
\end{enumerate}

\[
\vartheta (\exp x) = \sum_ {u \in \Gamma (S)} e ^ {- Q (x + u)}.
\]

\begin{enumerate}
\def\labelenumi{\alph{enumi})}
\setcounter{enumi}{2}
\tightlist
\item
  设 A 是 t 的一个单房，dx 是 t 上的一个 Haar 测度，h 是 t 上在仿射 Weyl 群 \(W_{a}^{\prime}\) 下不变的局部可积函数（§5，节 2）。证明等式
\end{enumerate}

\[
\int_ {\mathrm{A}} h (x) d x = \frac {1}{w (\mathrm{G})} \int_ {\mathfrak {t}} h (x) e ^ {- \mathrm{Q} (x)} (\vartheta_ {0} (x)) ^ {- 1} d x.
\]

\(d\) ) 对 \(x\in \mathfrak{g}\)，置 \(\xi (x) = \lambda_{\mathfrak{g}}(x)e^{-\mathrm{Q}(x)}(\vartheta (\exp x))^{-1}\)，其中 \(\lambda_{\mathfrak{g}}(x) = \det \frac{e^{\operatorname*{ad}x} - 1}{\operatorname*{ad}x}\)（§6，节 3）。证明 \(\xi\) 是一个 \(\mathrm{C}^\infty\) 类函数（\(\mathfrak{g}\) 上的），且若 \(\mu\) 是 \(\mathfrak{g}\) 上的一个 Haar 测度，则 \(\exp_{\mathrm{G}}\) 下测度 \(\xi \mu\) 的像是 G 上的一个 Haar 测度（利用 \(c\) ) 以及定理 1 的推论 2 与命题 4（§6，节 2 与节 3））。

\begin{enumerate}
\def\labelenumi{\alph{enumi})}
\setcounter{enumi}{4}
\tightlist
\item
  证明公式 \(\vartheta_0(x) = m\sum_{\lambda \in \mathrm{X(T)}}\exp (\delta (\lambda)(x) - \frac{1}{4}\mathbf{Q}'(\delta (\lambda)))\)（对 \(x\in \mathfrak{t}\)），其中 \(\mathbf{Q}'\) 是 \(\mathfrak{t}_{\mathbf{C}}^{*}\) 上的二次型，它是 \(\mathfrak{t}_{\mathbf{C}}\) 上由 \(\mathbf{Q}\) 诱导的二次型的逆型，\(m\) 是一个应当算出的常数（利用 Poisson 公式，参见~《谱理论》（编纂中）第 II 章，§1，节 8）。推出等式 \(\vartheta (t) = m\sum_{\lambda \in \mathrm{X(T)}}e^{-\mathrm{Q}'(\delta (\lambda)) / 4}t^{\lambda}\)（对 \(t\in \mathrm{T}\)）。
\end{enumerate}

\begin{enumerate}
\def\labelenumi{\arabic{enumi})}
\setcounter{enumi}{4}
\tightlist
\item
  \begin{enumerate}
  \def\labelenumii{\alph{enumii})}
  \tightlist
  \item
    设 V 是一个实向量空间，f 是 V 上的一个非零线性形式，H 是 f 的核。~那么，函数 \(\varphi\)（V 上的，\(C^{\infty}\) 类）在 H 上为零当且仅当它可写成 \(f\varphi'\)，其中 \(\varphi'\) 是 V 上的一个 \(C^{\infty}\) 类函数。
  \end{enumerate}
\end{enumerate}

\begin{enumerate}
\def\labelenumi{\alph{enumi})}
\setcounter{enumi}{1}
\item
  设 V 是一个实向量空间，\((f_{i})_{i\in\mathrm{I}}\) 是一个由非零线性形式组成的有限族，使得诸 \(H_{i}=\operatorname{Ker}f_{i}\) 两两不同。那么，函数 \(\varphi\)（V 上的，\(C^{\infty}\) 类）在诸 \(H_{i}\) 的并上为零当且仅当它可写成 \(\psi\prod_{i\in I}f_{i}\)，其中 \(\psi\) 是 V 上的一个 \(C^{\infty}\) 类函数。
\item
  设 T 是一个环面，\((\alpha_{i})_{i\in I}\) 是一个由 T 的异于 1 的特征标组成的有限族，使得诸 \(K_{i} = Ker \alpha_{i}\) 两两不同。那么，函数 \(\varphi\)（T 上的，\(C^{\infty}\) 类）在诸 \(K_{i}\) 的并上为零当且仅当它可写成 \(\psi \cdot \prod_{i\in I}(\alpha_{i}-1)\)，其中 \(\psi\) 是 T 上的一个 \(C^{\infty}\) 类函数（在 T 上局部地论证并利用 b)）。
\end{enumerate}

\(d\) ) 用节 4 的记号，证明映射 \(b_{\infty}:\mathcal{C}^{\infty}(\mathrm{T})^{\mathrm{W}}\to \mathcal{C}^{\infty}(\mathrm{T})^{-\mathrm{W}}\) 是双射。

\begin{enumerate}
\def\labelenumi{\arabic{enumi})}
\setcounter{enumi}{5}
\item
  假定 G 不是交换的。证明 T 上的连续函数 \(\mathrm{J}(\rho)^{1/3}\) 在 W 下是反不变的，但不属于映射 \(b_{c}: \mathcal{C}(\mathrm{T})^{\mathrm{W}} \to \mathcal{C}(\mathrm{T})^{-\mathrm{W}}\) 的像。
\end{enumerate}

\exercisesection{\S~9}

\begin{enumerate}
\def\labelenumi{\arabic{enumi})}
\item
  设 A 是 R 的由 0 与诸实数 \(1/n\)（\(n\) 为 \(\geq 1\) 的整数）组成的紧子集。证明当 \(\mathcal{C}^r(\mathbf{R};\mathbf{R})\) 赋予 A 上的一致 \(\mathrm{C}^r\) 收敛拓扑时，在 A 的邻域中为嵌入的态射组成的集合在 \(\mathcal{C}^r(\mathbf{R};\mathbf{R})\) 中不是开的（考察一个函数序列 \((f_n)_{n\geq 1}\)，使得 \(f_n(x) = x\)（对 \(x\leq \frac{1}{n + 1}\)），\(f_n(x) = x - \frac{1}{n}\)（对 \(x\geq \frac{1}{n}\)））。\\
\item
  设 X 是一个 \(\mathrm{C}^r\) 类（\(1\leq r\leq \infty\)）的分离流形，它无穷可数且是维数 \(n\) 的纯流形。
\end{enumerate}

\begin{enumerate}
\def\labelenumi{\alph{enumi})}
\item
  假定存在 X 到一个有限维向量空间 V 中的嵌入 \(\varphi\)。证明存在 X 在 \(\mathbf{R}^{2n + 1}\) 中的一个嵌入（若 \(\dim \mathrm{V} > 2n + 1\)，证明存在一点 \(p\)（V 中的），使得对所有 \(x \in \mathrm{X}\)，连接 \(p\) 与 \(\varphi(x)\) 的直线与 \(\varphi(\mathrm{V})\) 只交于点 \(\varphi(x)\)，且在该点横截相交；由此推出存在 X 在一个维数等于 \(\dim \mathrm{V} - 1\) 的空间中的嵌入）。
\item
  证明存在 X 在 \(R^{2n+1}\) 中的一个嵌入。（设 O 是 X 的开子集组成的集合，U 是 O 中由这样的开集 U 组成的子集：存在一个在 U 上的限制为嵌入的态射 \(\varphi: X \to R^{2n+1}\)；利用 a)，证明 U 是 O 的一个拟全子集（《一般拓扑学》第 IX 章，§4，习题 27），从而等于 O。）
\item
  证明存在 X 在 \(\mathbf{R}^{2n + 1}\) 中的一个真嵌入。（借助 X 上的一个真函数构造 X 在 \(\mathbf{R}^{2n + 2}\) 中的一个真嵌入。）
\end{enumerate}

¶ 3) 设 G 是一个李群，H 是 G 的一个紧子群。假定 G 具有一个忠实的（有限维）线性表示。

\begin{enumerate}
\def\labelenumi{\alph{enumi})}
\item
  设 \(\Theta_{\mathrm{H}}(\mathrm{G})\) 是 \(\mathcal{C}(\mathrm{H};\mathbf{R})\) 中由 G 上的（连续）代表函数在 H 上的限制组成的子代数。证明 \(\Theta_{\mathrm{H}}(\mathrm{G})\) 在 \(\mathcal{C}(\mathrm{H};\mathbf{R})\) 中对一致收敛拓扑是稠密的。
\item
  设 \(f \in \Theta_{\mathrm{H}}(\mathrm{G})\)。证明存在 G 在一个有限维实向量空间上的表示 \(\sigma\)，使得表示 \(\sigma|H\) 的某个系数与 f 不正交（《谱理论》（编纂中））。
\item
  设 \(\rho:H\to\mathbf{GL}(V)\) 是 H 在一个有限维实向量空间上的表示。证明存在一个有限维实向量空间 W、一个表示 \(\sigma:G\to\mathbf{GL}(W)\) 与一个单同态 \(u:V\to W\)，使得 \(u(\rho(h)v)=\sigma(h)u(v)\)（对 \(h\in H\)，\(v\in V\)）。
\end{enumerate}

（化归到 \(\rho\) 不可约的情形，并利用 b)。）

\begin{enumerate}
\def\labelenumi{\arabic{enumi})}
\setcounter{enumi}{3}
\item
  设 G 是一个李群，H 是 G 的一个紧子群，\(\rho\) 是 H 在一个实希尔伯特空间 V 上的一个酉表示。证明存在一个实希尔伯特空间 W、G 在 W 上的一个酉表示 \(\sigma\) 以及一个具有闭像的等距单射 \(u: V \to W\)，使得 \(u(\rho(h)v) = \sigma(h)u(v)\)（对 \(h \in H\)，\(v \in V\)）（方法与习题 3 相同）。
\item
  设 G 是一个李群，H 是 G 的一个紧子群。假定存在 G 在一个有限维实向量空间上的一个忠实线性表示 \(\rho: G \to \mathbf{GL}(V)\)。
\end{enumerate}

\begin{enumerate}
\def\labelenumi{\alph{enumi})}
\item
  证明存在表示 \(\sigma\)（\(\mathbf{GL}(\mathrm{V})\) 在一个有限维实向量空间 W 上的）、V 上一个分离的正二次型 q 以及 W 中的一个向量 w，使得 \(\rho(\mathrm{H}) = \mathbf{O}(q) \cap \mathrm{F}_{w}\)，其中 \(F_{w}\) 是 w 在 \(\mathbf{GL}(\mathrm{V})\) 中的稳定子群（选取一个在 H 下不变的二次型 q 以及 \(\mathbf{O}(q)\) 的一个使 \(\rho(\mathrm{H})\) 为一点的稳定子群的表示（节 2 的推论 2），然后应用习题 3 c)）。
\item
  由 a) 推出，存在一个有限维实向量空间 E、G 在 E 上的一个表示以及一个以 H 为稳定子群的向量 \(e \in E\)（取 \(E = W \oplus Q\)，其中 Q 是 V 上二次型组成的空间，且 \(e = (w, q)\)）。
\end{enumerate}

\begin{enumerate}
\def\labelenumi{\arabic{enumi})}
\setcounter{enumi}{5}
\item
  设 G 是一个李群，H 是 G 的一个紧子群。证明存在 G 在一个实希尔伯特空间 E 上的一个酉表示以及一个以 H 为稳定子群的向量 \(e \in E\)（取 \(\mathrm{E} = \mathrm{L}^{2}(\mathrm{G})\)）。
\item
  设 G 是一个李群，H 是 G 的一个紧子群，\(\rho\) 是 H 在一个实希尔伯特空间 W 上的一个酉表示。用 X 表示流形 \(G \times ^{H}W\)（节 3）。
\end{enumerate}

\begin{enumerate}
\def\labelenumi{\alph{enumi})}
\item
  证明存在一个希尔伯特空间 V、G 在 V 上的一个酉表示 \(\sigma\) 以及一个（解析）嵌入 \(\varphi: X \to V\)，使得 \(\varphi(gx) = \sigma(g)\varphi(x)\)（对 \(g \in G\)，\(x \in X\)）（利用习题 4 与习题 6）。
\item
  证明若 W 是有限维的且 G 具有一个忠实的有限维线性表示，则 V 可取为有限维的（利用习题 3 与习题 5）。
\end{enumerate}

¶ 8) 设 X 是一个 \(C^{r}\) 类（\(1 \leq r \leq \infty\)）的仿紧流形，G 是一个真地作用在 X 上的李群，且作用律 \((g, x) \mapsto gx\) 是 \(C^{r}\) 类的。

\begin{enumerate}
\def\labelenumi{\alph{enumi})}
\item
  证明存在 G 在一个希尔伯特空间 V 上的一个酉表示 \(\rho\) 以及一个嵌入 \(\varphi\)（\(C^{r}\) 类的，X 到 V 中的），使得 \(\varphi(gx) = \sigma(g)\varphi(x)\)（对 \(g \in G, x \in X\)）。（利用命题 6、习题 7，并按命题 4 的证明方式论证。）
\item
  再作如下补充假定：
\end{enumerate}

\begin{enumerate}
\def\labelenumi{(\roman{enumi})}
\item
  G 具有一个有限维线性表示；
\item
  \(\mathrm{X}\) 无穷可数且维数有界；
\item
  X 在 G 的作用下只有有限多个轨道类型。
\end{enumerate}

证明存在 X 的一个由在 G 下稳定的开集 \((\mathrm{U}_{i})_{i\in\mathrm{I}}\) 组成的有限覆盖、G 的紧子群 \((\mathrm{H}_{i})_{i\in\mathrm{I}}\)，以及对每个 \(i\in I\)，一个闭子流形 \(S_{i}\)（\(U_{i}\) 的，在 \(H_{i}\) 下稳定），使得映射 \((g,s)\mapsto gs\) 通过过渡到商空间诱导出从 \(G\times^{H_{i}}S_{i}\) 到 \(U_{i}\) 的一个同构（证明 X/G 中其原像在 X 内容许这样一个覆盖的开子集组成 X/G 的开子集集合的一个拟全子集，参见~《一般拓扑学》第 IX 章，§4，习题 27）。

\begin{enumerate}
\def\labelenumi{\alph{enumi})}
\setcounter{enumi}{2}
\tightlist
\item
  在 b) 的假定下，证明存在 G 在一个有限维实向量空间 V 上的一个表示以及 X 在 V 中的一个与 G 的作用相容的 \(C^{r}\) 类嵌入（在 G 的子群组成的集合上用归纳法论证，利用 b)、习题 2 与习题 7）。
\end{enumerate}

\begin{enumerate}
\def\labelenumi{\arabic{enumi})}
\setcounter{enumi}{8}
\tightlist
\item
  设 \(G\) 是一个真地作用在流形 \(X\) 上的李群。作下列两个假定之一：
\end{enumerate}

\begin{enumerate}
\def\labelenumi{(\roman{enumi})}
\tightlist
\item
  X/G 是紧的，且 X 无穷可数、维数有界；\\
\item
  X 是一个 G 在其上线性作用的有限维向量空间。
\end{enumerate}

证明 X 的元素的轨道类型组成的集合是有限的（同时对两种情形按 dim X 用归纳法处理，并利用命题 6）。

\[
\mathbf {G L} (3, \mathbf {R})
\]

\(\left(\begin{array}{ccc}1 & a & b \\ 0 & 1 & 0 \\ 0 & 0 & 1\end{array}\right), a, b \in \mathbf{R}.\) 证明，对 \(G\) 在

\(\mathbf{R}^3\) 上的恒等表示，轨道类型的个数是无限的。

\end{enumerate}
% semantic-fragments: legacy-input-seam
\relax{}
¶ 11) 对任意整数 \(n \geq 1\) ，定义映射 \(\varphi_n\)，其定义由 \((0, \frac{1}{2} (\times \mathbf{T}^2 \text{ 到 } \mathbf{R}^3 \text{ 令 } \varphi_n(r; \alpha, \beta) = ((n + r \cos 2\pi\beta) \cos 2\pi\alpha, (n + r \cos 2\pi\beta) \sin 2\pi\alpha, r \sin 2\pi\beta)\) 给出。令

\[
\mathrm{T} _ {n} = \varphi_ {n} \left(\left[ 0, \frac {1}{2} \right[ \times \mathbf {T} ^ {2}\right), \quad \mathrm{S} _ {n} = \varphi_ {n} (\{0 \} \times \mathbf {T} ^ {2}).
\]

\begin{enumerate}
\def\labelenumi{\alph{enumi})}
\item
  证明 \(\varphi_{n}\) 在 \(\left[0,\frac{1}{2}\right[\times T^{2}\) 上的限制是一个（类为 \(C^{\omega}\) 的）从 \(\left]0,\frac{1}{2}\right[\times T^{2}\) 到 \(T_{n}-S_{n}\) 的同构。
\item
  设 \(f_{n}:\mathrm{T}_{n}\to \mathrm{T}_{n}\) 为在 \(\mathrm{S}_n\) 上与恒等映射重合并满足下式的映射：
\end{enumerate}

\[
f _ {n} (\varphi_ {n} (r; \alpha , \beta)) = \varphi_ {n} (r; \alpha - (n - 1) \beta , - \alpha + n \beta)
\]

当 r \textgreater{} 0 时，证明 \(f_{n}\) 诱导出 \(T_{n} - S_{n}\) 的一个自同构。

\begin{enumerate}
\def\labelenumi{\alph{enumi})}
\setcounter{enumi}{2}
\tightlist
\item
  证明在 \(R^{3}\) 上存在一个实解析流形结构，使得映射 \(f_{n}: T_{n} \to R^{3}\) 和典则注入 \(R^{3}-\bigcup_{n} S_{n} \to R^{3}\) 都是解析的。将由此定义的实解析流形记为 X。
\end{enumerate}

\(d)\) 对于 \(\vartheta \in \mathbf{T}\)，将旋转记为 \(\mathrm{R}_{\vartheta}\)

\[
(x, y, z) \mapsto (x \cos 2 \pi \vartheta - y \sin 2 \pi \vartheta , x \sin 2 \pi \vartheta + y \cos 2 \pi \vartheta , z)
\]

在 \(\mathbf{R}^3\) 上。证明，对于 \(\theta \in \mathbf{T}\) 和 \(u \in \mathrm{X}\)，规定

\[
\vartheta . u = \mathrm{R} _ {\vartheta} (u) \quad \text { if } \quad u \in \mathrm{X} - \bigcup_ {n} \mathrm{S} _ {n};
\]

\(\vartheta .u = \mathrm{R}_{n\vartheta}(u)\) 若 \(u\in \mathrm{S}_n\) 对某个整数 \(n\geq 1\)

定义了 \(\mathbf{T}\) 在 \(\mathrm{X}\) 上的解析作用律。

\begin{enumerate}
\def\labelenumi{\alph{enumi})}
\setcounter{enumi}{4}
\item
  证明 X 的轨道类型集合（对于 d) 中定义的 T 作用）是无限的。
\item
  证明不存在一个与 T 的作用相容的嵌入，将 X 嵌入某个 T 在线性作用的有限维向量空间中（使用习题 9）。
\end{enumerate}

\begin{enumerate}
\def\labelenumi{\arabic{enumi})}
\setcounter{enumi}{11}
\tightlist
\item
  设 \(\mathrm{G}\) 为紧李群。
\end{enumerate}

\begin{enumerate}
\def\labelenumi{\alph{enumi})}
\item
  证明整子群的正规化子的共轭类集合是有限的（考察 G 在 L(G) 的子空间 Grassmann 流形上的作用，并应用习题 9）。
\item
  （紧）半单子群的共轭类集合是有限的（注意到一个李代数至多包含有限多个半单理想（第一章，§6，习题 7），并使用 a)）。
\end{enumerate}

\begin{enumerate}
\def\labelenumi{\arabic{enumi})}
\setcounter{enumi}{12}
\item
  设 G 为李群，H、K 为 G 的两个紧子群。假设 G 具有忠实的有限维线性表示。证明，存在 H 的子群的有限集合 F，使得对于所有 \(g \in G\)，子群 \(H \cap gKg^{-1}\) 在 H 中共轭于 F 中的某个子群（使用习题 5 和 9）。
\item
  假设流形 X 是仿紧的且局部有限维。设 G 为在 X 上真作用的李群，H 为 G 的紧子群，t 为 H 的共轭类。
\end{enumerate}

\begin{enumerate}
\def\labelenumi{\alph{enumi})}
\item
  证明 X 中固定子群等于 H 的点所成的集合 \(X_{H}\) 是 X 的局部闭子流形。
\item
  证明映射 \((g,x)\mapsto gx\) 从 \(G\times X_{H}\) 到 X 诱导出一个（类为 \(C^{r}\) 的）从 \(G\times^{\mathrm{N(H)}}X_{\mathrm{H}}\) 到 \(\mathrm{X}_{(t)}\) 的同构。
\end{enumerate}

¶ 15) 设 G 为在拓扑空间 E 上真作用的局部紧拓扑群；设 \(\rho\) 是 G 在有限维实向量空间 V 上的一个表示。记 G 上的右 Haar 测度为 dg。

\begin{enumerate}
\def\labelenumi{\alph{enumi})}
\tightlist
\item
  设 \(\mathcal{P}\) 为 E 的具有如下性质的子集 A 所成的集合：存在 E 的一个开覆盖 \((\mathrm{U}_{\alpha})_{\alpha \in \mathrm{I}}\)，使得对于所有 \(\alpha \in \mathrm{I}\)，\(g \in \mathrm{G}\) 中满足 \(g\mathrm{A} \cap \mathrm{U}_{\alpha} \neq \emptyset\) 的元素所成的集合在 G 中相对紧。
\end{enumerate}

证明，对于每个支集属于 P 的连续函数 \(f: E \to V\)，映射 \(x \mapsto \int_{G} \rho(g)^{-1} f(gx) dg\) 是从 E 到 V 的连续映射，并且与 G 的作用相容。

\begin{enumerate}
\def\labelenumi{\alph{enumi})}
\setcounter{enumi}{1}
\tightlist
\item
  作如下两种假设之一：
\end{enumerate}

\begin{enumerate}
\def\labelenumi{(\roman{enumi})}
\item
  空间 E/G 是正则的；
\item
  存在 E 的一个开子集 U 和 G 的一个紧子集 K，使得 \(\mathrm{E} = \mathrm{GU}\)，并且 \(g\mathrm{U}\cap \mathrm{U} = \varnothing\) 对 \(g\notin \mathrm{K}\) 成立。
\end{enumerate}

证明 E 中每一点 x 都有一个属于 P 的邻域（在情形 (i) 中，取 \(A = V \cap W\)，其中 V 是 x 的一个邻域，当 g 在 G 的某个紧子集之外时满足 \(gV \cap V = \varnothing\)；W 取为 Gx 的一个在 G 下稳定且包含于 GV 中的闭邻域）。

E 中每一点都有一个在 G 下稳定且满足 (ii) 的开邻域。

\begin{enumerate}
\def\labelenumi{\alph{enumi})}
\setcounter{enumi}{2}
\tightlist
\item
  另外假设 E 是完全正则的。设 \(x \in E, v \in V\)，使得 x 的固定子群包含于 v 的固定子群。证明存在连续映射 \(F : E \to V\)，它与 G 的作用相容，且 \(\mathrm{F}(x) = v\)。（令 F 为 E 上支集属于 P 的连续数值函数空间，并令 \(u : F \to V\) 为映射 \(\alpha \mapsto \int_{G} \alpha(gx) \rho(g^{-1}). v \, dg\)。设 C 是 V 中 v 的一个凸邻域；构造一个属于 P 的 x 的邻域 A，使得 \(gx \in A\) 蕴含 \(\rho(g^{-1}). v \in C\)，并构造 E 上的一个函数 \(\alpha\)，其支集包含于 A，满足 \(\alpha(x) \neq 0\) 且 \(\int_{G} \alpha(gx) \, dg = 1\)。证明 \(u(\alpha)\) 属于 C，并推出 \(v \in Im u\)。）
\end{enumerate}

\begin{enumerate}
\def\labelenumi{\arabic{enumi})}
\setcounter{enumi}{15}
\tightlist
\item
  设 G 为在分离拓扑空间 E 上真作用的拓扑群。设 x 是 E 中一点，H 是它在 G 中的固定子群，S 是包含 x 且在 H 下稳定的 E 的子集。群 H 右作用于 \(G \times S\)，作用公式为 \((g, s).h = (gh, h^{-1}s)\)，其中 \(g \in G, h \in H, s \in S\)；映射 \((g, s) \mapsto gs\) 通过取商诱导出从 \((\mathrm{G} \times \mathrm{S})/\mathrm{H}\) 到 X 的映射。若此映射是到 X 的某个开子集上的同胚，则称 S 是 x 处的横截集。
\end{enumerate}

\begin{enumerate}
\def\labelenumi{\alph{enumi})}
\item
  证明，若 \(\mathbf{G}\) 是离散的，则 \(x\) 处存在一个横截集。
\item
  设 F 是一个分离拓扑空间，G 在其上真作用；设 \(f : E \to F\) 是与 G 的作用相容的连续映射，且 \(f(x)\) 在 G 中的固定子群等于 H。证明，若 S 是 \(f(x)\) 处的横截集，则 \(f^{-1}(S)\) 是 x 处的横截集。
\item
  设 N 是 G 的闭正规子群，\(\pi: E \to E/N\) 是典则投影，T 是 E/N 上 G/N 作用下 \(\pi(x)\) 处的横截集，且 \(S \subset \pi^{-1}(T)\) 是 \(\pi^{-1}(T)\) 上 HN 作用下 x 处的横截集。证明 S 是 E 中 x 处的横截集（对于 G 的作用）。
\end{enumerate}

\begin{enumerate}
\def\labelenumi{\arabic{enumi})}
\setcounter{enumi}{16}
\tightlist
\item
  设 \(G\) 为李群。我们将证明 \(G\) 具有如下性质：
\end{enumerate}

\begin{enumerate}
\def\labelenumi{(\Alph{enumi})}
\setcounter{enumi}{19}
\tightlist
\item
  对于每个完全正则拓扑空间 \(\mathrm{E}\)，在其上 \(\mathrm{G}\) 真作用，并且对于每一点 \(x\)，在 \(\mathrm{E}\) 中的 \(x\) 处都存在一个横截集（习题 16）。
\end{enumerate}

\begin{enumerate}
\def\labelenumi{\alph{enumi})}
\item
  证明每个具有忠实有限维线性表示的李群都满足 (T)（使用习题 15 和 16 b)，以及命题 6）。
\item
  证明，若 G 有一个闭正规子群 N，使得 G/N 满足 (T)，并且对 G 的每个紧子群 K，KN 满足 (T)，则 G 满足 (T)（应用习题 16 c)）。
\item
  若 \(\mathrm{G}_0\) 是紧的，则 G 满足 (T)。
\item
  证明，若 G 有一个离散正规子群 N，使得 G/N 满足 (T)，则 G 满足 (T)。
\item
  证明 G 满足 (T)（令 N 为伴随表示的核；证明 \(G/N_{0}\) 满足 (T)，然后应用 a)、b) 以及 §1 的习题 9）。
\end{enumerate}

\begin{enumerate}
\def\labelenumi{\arabic{enumi})}
\setcounter{enumi}{17}
\tightlist
\item
  设 \(G\) 为在完全正则拓扑空间 \(E\) 上真作用的李群。
\end{enumerate}

\begin{enumerate}
\def\labelenumi{\alph{enumi})}
\item
  设 \(x \in E\)，且 t 是它的轨道类型；证明存在 x 的一个在 G 下稳定的开邻域 U，使得对于所有 \(u \in U\)，u 的类型都满足 \(\geq t\)。
\item
  假设 G 在 E 上自由作用；令 \(\pi: E \to E/G\) 为典则投影。证明，对于 E/G 中的每一点 z，都存在 z 的一个开邻域 U 和连续映射 \(s: U \to E\)，使得 \(\pi \circ s(u) = u\) 对所有 \(u \in U\) 都成立。
\end{enumerate}

\begin{enumerate}
\def\labelenumi{\arabic{enumi})}
\setcounter{enumi}{18}
\tightlist
\item
  设 G 为李群，H 为 G 的紧子群。证明存在 H 的一个邻域 V，使得 G 中包含于 V 的每个子群都共轭于 H 的某个子群（将习题 18 a) 应用于 G 的紧子集空间，参见~《积分》，第八章，§5，第 6 节）。
\end{enumerate}

¶ 20) 设 G 为紧李群，m 为正整数。

\begin{enumerate}
\def\labelenumi{\alph{enumi})}
\item
  证明 G 中满足阶 \(\leq m\) 的子群的共轭类集合是有限的（假设结论不成立，构造一个有限群 F 以及一列同态 \(\varphi_{n}: F \to G\)，使得 \(\varphi_{i}(F)\) 不共轭于 \(\varphi_{j}(F)\)（当 \(i \neq j\) 时），并且 \(\varphi_{n}(f)\) 都趋于极限 \(\varphi(f)\)，对所有 \(f \in F\) 均如此；证明 \(\varphi\) 是同态，并证明这将导致与习题 19 的矛盾）。
\item
  证明 G 中所有元素的阶均满足 \(\leq m\) 的子群 F 的共轭类集合是有限的（令 \(\mu\) 为 G 上的 Haar 测度，令 U 为单位元的一个对称邻域，使得 \(U^{2}\) 不包含阶满足 \(\leq m\) 的非平凡元素；证明 \(\operatorname{Card}(F) \leq \mu(G)/\mu(U)\)）。
\end{enumerate}

¶ 21) 设 G 为紧李群，T 为 G 的极大环面。

\begin{enumerate}
\def\labelenumi{\alph{enumi})}
\item
  设 S 是 G 的闭子群的一个在共轭下稳定的集合，并且子群族 \((\mathrm{S} \cap \mathrm{T})_{\mathrm{S} \in \mathcal{S}}\) 是有限的。证明子群 \(S_{0}\)（其中 \(S \in S\)）的共轭类集合是有限的（利用习题 12 a)），化归到子群 \(S_{0}\) 都是正规的情形；考察群 \(\mathrm{C}(S_{0})_{0}\) 和 \(\mathrm{D}(S_{0})\)，并应用习题 12 b)）。
\item
  证明 S 是 G 的子群的有限个共轭类的并（利用 a)，化归到子群 \(S \in S\) 都具有同一个单位元连通分支 \(\Sigma\) 的情形，其中该分支在 G 中正规；然后给出群 \(S/\Sigma\) 中元素的阶的上界，并应用习题 20 b)）。
\item
  设 E 是一个 G 在其上连续作用的分离拓扑空间。证明，若 E 的点关于 T 的作用只有有限种轨道类型，则关于 G 的作用也只有有限种轨道类型。
\end{enumerate}

